\section{B. Metabelian torsion and full peripheral lifts}

\begin{statement}{B.1}{Open}{Metabelian classification and the remaining arithmetic centralizer problem}
The metabelian pure peripheral group, its symmetric multiplier-one
subgroup, and their finite-order units are described in B.7--B.10.
This proves the metabelian assertions in Leads 5, 6, 7 and the unit
assertion in Lead 22, with the nonzero condition in Lead 7 made
explicit. The R28 arithmetic families and all its projected
rigidity assertions are independently verified in B.18 and
B.20--B.24. The least derived level at which the centralizer
itself is geometric remains unresolved.
\end{statement}
\begin{proof}[Scope of the result and the unfinished step]
The unit descriptions are exact statements about $M$. Their
full multiplier-one peripheral image is separately determined
by B.17, and B.18 classifies the full lifts in both R28 families.
Commuting with $G_K$ is an additional condition: B.20 excludes
every nonidentity family member already at level three.
The independent arithmetic-cover arguments of B.20--B.22 prove
that the level-three centralizer projects to $\mathfrak S_3$
on $M$, and the level-four centralizer projects to
$\mathfrak S_3$ on $S_3$. The same-level peripheral kernel of
level three to level two is still not computed. B.23 explains
why the explicit nontrivial unrestricted kernel classes do not
belong to it. A next attack must either construct a peripheral
last-layer cocycle surviving arithmetic or show that the three
peripheral constraints kill all such cocycles; neither follows
from vanishing of the projected image.
\end{proof}



\begin{statement}{B.2}{Open}{Beta identities versus the complete metabelian pentagon image}
For multiplier one, the metabelian two-cycle and three-cycle
identities are satisfied by $\varphi_0$ and by every $\varphi_g$
in B.10. The beta two-cocycle identity forced by the pentagon fails
for $\varphi_0$ and holds for every such $\varphi_g$.
Consequently $\varphi_0$ is outside the metabelian image of
$\widehat{\mathrm{GT}}$. This does not by itself prove that any
$\varphi_g$ belongs to that image. The equivalence of the listed
beta identities with the complete requested metabelian pentagon
condition is not established here.
\end{statement}
\begin{proof}[Normalization and source scope]
The action on $M'$ in this paper is the unit $B'_F$ of
\cite[\S2.3, p.~101]{PBC_Furusho2003}. Its normalization there is
\[
 B'_F=S_\lambda(X)S_\lambda(Y)B_F,
 \qquad S_a(T)=\frac{T^a-1}{T-1}.
\]
Thus it agrees with the adelic beta function $B_F$ at multiplier
one. The same opened passage gives the beta--gamma factorization
and identifies the original statement as Ihara, Proposition 1.6.1.
Cancellation in that factorization proves
\[
 B_F(X,Y)B_F(XY,W)=B_F(Y,W)B_F(X,YW).
\]
The character realization is injective, so this is an identity in
the completed three-variable group algebra. The original Ihara
Part I text has not been opened; our verified source for this
specific theorem and its normalization is the research paper just
cited. Its statement is distinguished from a claim that the beta
identity alone reconstructs a full pentagon pair. B.9 supplies the
explicit failure, and B.10 proves the coboundary identity directly.
\end{proof}

\begin{statement}{B.3}{Proved}{A full lift and an explicit finite obstruction to an involutive lift}
The involution $\varphi_0$ of R7 lifts to the full profinite
gentle group $\widehat{\mathrm{GT}}_0\cap\mathcal P_1$
by B.27, necessarily with infinite outer
order. Its exact central obstruction is zero. More generally,
B.17 computes the obstruction for every multiplier-one
metabelian peripheral unit, and B.18 determines it for every
member of the two R28 families. The explicitly presented
quotient of order $614400$ in B.19 excludes the entire prescribed
involutive system (A--I), with forced descent of every hypothetical
full solution proved there. Thus Task B.3 is settled in both its
unrestricted-lift and finite-order formulations.
\end{statement}
\begin{proof}
B.12 gives an independent finite-order unitary lifting argument;
B.17 gives the more general exact formula, and B.11 supplies the
full source lift from the universal central cover. B.13 excludes
finite-order full lifts. B.19 provides the specified finite
obstruction and its executed certificate. These statements use
full profinite groups, with the mixed-prime quotient retained.

The absolute pronilpotent lifting theorem P.3 alone would not
prove this result: $M$ has nonnilpotent finite quotients, such as
$Q_{75}$, so there is no quotient map from
$F/R_{\mathrm{nil}}(F)$ onto $M$. The central criterion of B.11
and the integral computation of B.17 are the required replacement.
Consequently the tripod has no nonidentity finite pure peripheral
action of multiplier one, even though the full image on its
metabelian quotient contains the displayed nontrivial involutions.
\end{proof}

\begin{statement}{B.4}{Open}{Parity and the conjugacy classes of involutions}
R25 proves that every nonidentity finite-order element of
$\mathcal P$, $\widehat{\mathrm{GT}}_0$, or
$\widehat{\mathrm{GT}}$ is an involution of global multiplier $-1$.
It does not settle whether every gentle involution has $\Delta=0$,
or whether every involution of either Grothendieck--Teichm\"uller
group is conjugate to complex conjugation. These conclusions remain
unresolved here.
\end{statement}
\begin{proof}[The precise distinction]
For $\tau=(-1,f_\tau)$ the checked R25 formula is
$\Delta(\tau)=\delta((1,f_\tau))$. Conjugacy to $c$ implies
zero because $\Delta$ is a character and $\Delta(c)=0$.
The converse is not a consequence of a character computation.
Lead 23 needs an equivariant lift of a relative prime-component
construction through a kernel on which the order-two cohomology
is controlled. B.11 lifts automorphisms and conjugators, but does
not lift a finite group action. In particular it does not imply
that a chosen lift satisfies $\sigma_2c\sigma_2c=1$.
B.29 sharpens the target: $\tau_0=\varphi_0c$ is a negative
metabelian involution in the full gentle image which is not
conjugate to $c$ in the image of $\mathcal P$. Its actual
full involutive lifting remains open. B.28 proves that its
order-$75$ residue is Galois-realized, so no test retaining
only that residue, even over an entire relative pro-two
tower, can decide this lifting question. A new obstruction
must retain at least one further metabelian block, such as
the order-$100$ quotient used in B.29.
\end{proof}

\begin{statement}{B.5}{Open}{The centralizer is known but the full conjugacy problem remains}
The centralizer of complex conjugation in $\mathcal P$ is
$\langle c\rangle$, as proved in B.14. The specific four-puncture
involution of R22 does not descend by filling one puncture.
The generation criterion in the task is proved in B.15.
The asserted existence and conjugacy of the glued involution of
Lead 23, the set $H^1(\langle c\rangle,\mathcal H)$, and its
complete compatible-family analysis at the indicated relative
levels remain unresolved.
\end{statement}
\begin{proof}[Where the reductions stop]
The endomorphism $x\mapsto x$, $y\mapsto hyh^{-1}$ constructed
in R25 is an automorphism of every maximal pro-prime quotient.
This does not imply that its image generates every mixed-prime
finite quotient. The third peripheral equation and full generation
are exactly the missing conditions. P.3 can lift a prescribed
pronilpotent peripheral map, but does not require that the lift
conjugate a specified full involution. At the cohomological level,
the classes of multiplier-$-1$ involutions modulo conjugation by
$\mathcal H$ are the nonabelian cohomology set
$H^1(\langle c\rangle,\mathcal H)$. The character induced by
$\Delta$ is constant on these classes. Neither its surjectivity
nor its fibres follow from surjectivity of the underlying group
homomorphism on a relative quotient.
B.25 proves the inverse-system and odd-prime conjugacy
reductions in Lead 26(b). B.26 refutes the suggestion that
$\Delta$ is already a well-defined invariant of an unframed
element at the specified first relative level.
\end{proof}

\begin{statement}{B.6}{Open}{Ghost values require the actual composition law}
The ghost structure theorem in Task B.6 is not established here.
In particular the asserted equivalence between abelianness of a
ghost quotient and abelianness of three ambient centralizer
quotients has not been proved.
\end{statement}
\begin{proof}[Exact algebra that must be retained]
For actual peripheral automorphisms $\Phi,\Psi$ with chosen frames
$b_i(\Phi),b_i(\Psi)$, the frame of their composite is
\[
 b_i(\Phi\Psi)=\Phi(b_i(\Psi))\,b_i(\Phi),
\]
up to the indicated right peripheral centralizer, with a fixed
choice of inverse convention. Thus the values form a crossed
composition law. Componentwise multiplication of ambient
centralizers is not this law unless the action on those values
has also been shown trivial. Equality of the three generator
images determines the induced automorphism at a retained source
level, but an arbitrary triple of centralizer elements need not
have the common base values required for a fibre-product lift.
R24 E.58 proves realization when those source and base conditions
hold. A general injectivity and image theorem modulo the proposed
$V$-ghosts needs an exact definition of their chosen frame quotient
and this crossed law. No equivalence is inferred merely from the
centralizers being abelian or nonabelian.
\end{proof}

\subsection{The exact metabelian unit calculation}

\begin{statement}{B.7}{Proved}{Pure peripheral automorphisms as semilinear units}
Write $R=\widehat{\mathbb Z}[[\widehat{\mathbb Z}^2]]$,
$d=(X-1)(Y-1)(XY-1)$, and $c_0=[x,y]$.
The multiplier-one pure peripheral outer group of $M$ is exactly
\[
                 \mathcal U=(1+dR)\cap R^\times,
\]
with multiplication of units. Its centralizer of the geometric
$S_3$ is $\mathcal U^{S_3}$, where $S_3$ permutes
$X,Y,(XY)^{-1}$. The whole pure peripheral group is the following
explicit group of pairs $(\lambda,\beta)$.
Put $\nu=(\lambda-1)/2$ and
\[
 k_\lambda=\frac{S_\lambda(X)-S_\lambda(XY)}{Y-1}
                     +X S_\nu(X)S_\lambda(XY).
\]
There must be an $r\in R$ such that
\[
 \begin{split}
 \beta&=S_\lambda(X)S_\lambda(Y)
                  +(X^\lambda-1)(Y^\lambda-1)r\in R^\times,\\
 X^\lambda(Y^\lambda-1)r-k_\lambda&\in (XY-1)R.
 \end{split}
\]
Its product is
$(\lambda,\beta)(\mu,\gamma)
=(\lambda\mu,\beta\,\gamma(X^\lambda,Y^\lambda))$.
Its symmetric subgroup is characterized by
$\beta^\sigma/\beta\in\langle X,Y\rangle$ for $\sigma\in S_3$.
Here $\langle X,Y\rangle$ denotes the closed monomial subgroup
$\widehat{\mathbb Z}^2$ of $R^\times$.
\end{statement}
\begin{proof}
We use the free rank-one $R$-module structure on $M'$ verified in
\cite[Corollary 4.1.5]{PBC_ChenDeligne2017}. The factors $X-1$,
$Y-1$, and $XY-1$ are non-zero-divisors. This follows either from
the regular completed-group-algebra sequence or from the product
of unramified power-series rings described in B.8.

The centralizer of each primitive element $x,y,xy$ in $M$ is its
procyclic subgroup. For example write an element uniquely as
$x^a y^b c_0^r$. Commuting with $x$, followed by augmentation of
the commutator exponent, forces $b=0$; the remaining equation is
$(X-1)r=0$, so $r=0$. The cases $y,xy$ follow by changing free
basis. Thus an outer peripheral class has a unique representative
\[
 x\longmapsto x^\lambda,\qquad
 y\longmapsto (c_0^r)^{-1}y^\lambda c_0^r,
 \qquad r\in R.
\]
To see existence of the normalization, first fix the image of
$x$, then remove the two abelian exponents in the $y$-conjugator
by its left $y$-power and the remaining inner $x$-power. The same
centralizer calculation proves uniqueness. These normalized
representatives form a group section for every multiplier:
composition and inversion preserve $x\mapsto x^\lambda$ and
the condition that the $y$-conjugator belongs to $M'$. Thus
an equality of their outer classes is an equality of the
normalized automorphisms themselves.

Commutator collection gives the displayed formula for $\beta$.
The map is an automorphism precisely when $\beta$ is a unit,
since it is already an automorphism of the quotient $A$ and is
$\lambda$-semilinear on its kernel $M'$. For positive integer $a$,
collection gives
\[
 (xy)^a=c_0^{\kappa_a}x^a y^a,
 \qquad \kappa_a=\frac{S_a(X)-S_a(XY)}{Y-1}.
\]
The formula extends to $a\in\widehat{\mathbb Z}$ by continuity.
Its division is legitimate because the numerator vanishes at $Y=1$
and $Y-1$ is regular. The class-two quotient shows that a conjugator
for $(xy)^\lambda$ may, after multiplication by a power of $xy$,
be written $x^{-\nu}c_0^s$. Conjugating by it contributes
\[
 X S_\nu(X)S_\lambda(XY)+((XY)^\lambda-1)s
\]
to the exponent relative to $(xy)^\lambda$. Comparison with the
image $x^\lambda(c_0^r)^{-1}y^\lambda c_0^r$ proves precisely the
second condition in the statement. $S_\lambda(T)$ is a unit for
unit $\lambda$: on a nonprincipal prime-to-$p$ character block it
is a ratio of two units, and on the principal block its constant
term is $\lambda_p$. Consequently $((XY)^\lambda-1)R=(XY-1)R$.

When $\lambda=1$ the condition reduces to
$X(Y-1)r\in(XY-1)R$. In $R/(XY-1)$ the factor $Y-1$ is regular,
so $r\in(XY-1)R$. This yields exactly $1+dR$. The action on $M'$
shows the product formula. For any geometric permutation the
semilinear factor on $c_0$ is a signed monomial; comparing
composites modulo inner automorphisms therefore gives exactly
$\beta^\sigma/\beta\in A$. At multiplier one both units lie in
$1+dR$, whose intersection with $A$ is $\{1\}$, by evaluation
at $X=1$ and $Y=1$. This gives the asserted invariance condition.
\end{proof}

\begin{statement}{B.8}{Proved}{All finite-order multiplier-one peripheral units}
For each prime $p$ let $\chi=(a,b)$ range over pairs of roots of
unity of order prime to $p$, modulo simultaneous Frobenius.
Let $\mathcal O_\chi$ be the unramified coefficient ring generated
by $a,b$. All elements of $\mathcal T$ are exactly the units whose
$\chi$-block is a constant root of unity
$\zeta_{p,\chi}\in\mathcal O_\chi^\times$, subject to
\[
 \zeta_{p,\chi}=1\quad\text{if }a=1,\ b=1,
                          \text{ or }ab=1,
\]
and to one uniform finite bound on the orders of all these roots.
The subset $\mathcal T_0$ consists exactly of the choices invariant
under the permutations of $(a,b,(ab)^{-1})$, with Frobenius
covariance. This includes the prime two without discarding its
constant root $-1$.
\end{statement}
\begin{proof}
The usual prime-to-$p$ character idempotents give the topological
product decomposition
\[
 R=\prod_p\prod_{[\chi]}
               \mathcal O_\chi[[T_1,T_2]],
 \quad X=a(1+T_1),\quad Y=b(1+T_2).
\]
It is obtained first for finite prime-to-$p$ abelian quotients
by the separable roots of unity and then by inverse limit; the
pro-$p$ factor is the two-variable power-series algebra.
Each factor is a characteristic-zero domain. A finite-order unit
in such a factor is constant: if its constant term is $\zeta$,
the first nonzero positive-degree term of its $n$-th power is
$n\zeta^{n-1}$ times that term, and cannot vanish. Thus all such
units are constant roots of unity. They have a common finite order
exactly when their component orders have a common finite bound.

The element $d$ is a unit exactly when none of $a,b,ab$ is one.
On every other block it has a nonconstant regular factor. A constant
$\zeta-1$ divisible by that factor must be zero, as evaluation on
its zero locus shows. B.7 now gives exactly the conditions in the
statement. Permuting the three character values gives the
$S_3$ criterion, including the coefficient Frobenius action.
\end{proof}

\begin{statement}{B.9}{Proved}{The cubic involution is symmetric and fails the beta cocycle}
The idempotent $e_0$ of R7 is supported on
$(a,b,(ab)^{-1})=(\omega,\omega,\omega)$ and its inverse.
Hence $\beta_0=1-2e_0$ belongs to $\mathcal T_0$.
It fails the beta two-cocycle identity, and therefore no full
pentagon pair has metabelian image $\varphi_0$.
The unit formula in Lead 22 is correct:
\[
 \ker(\Out(M)\to\GL_2(\widehat{\mathbb Z}))
              \simeq ((1+I)\cap R^\times)/A.
\]
The three special directions in degree five map injectively to the
principal metabelian block. Their integral image in the normalized
commutator coordinate is
$(u+v)\langle u^2,2uv,v^2\rangle_{\mathbb Z}$.
In particular it is saturated away from two; at two it has index
two in $(u+v)\mathbb Z[u,v]_2$.
\end{statement}
\begin{proof}
The support description proves the symmetry. If $f=c_0^r$, the
metabelian two-cycle is $r(X,Y)=r(Y,X)$, and the three-cycle is
\[
 Y^{-1}r(Z,X)+Zr(Y,Z)+r(X,Y)=0.
\]
Substitution of $r=(\beta-1)/((X-1)(Y-1))$ and the symmetry of
$\beta$ prove both identities, since
$Y^{-1}(Y-1)+Z(X-1)+(Z-1)=0$. Thus this conclusion concerns the
actual metabelian cycle equations, not only a dimension count.
At $(X,Y,W)=(\omega,\omega,-1)$ the two sides of the beta cocycle
are respectively $-1$ and $1$. The independent script
\texttt{pbc\_meta.py} checks this evaluation and all $216$ triples
of sixth-root characters; it finds $24$ failures for $\beta_0$.
B.2 supplies the necessary implication from the pentagon.

For the general IA group, an IA endomorphism sends the two
generators to $c_0^{r_1}x,c_0^{r_2}y$ and has determinant
$1+r_1(1-Y)+r_2(X-1)$. It is an automorphism exactly when this is
a unit, and is inner exactly when the determinant is a monomial.
These are \cite[Propositions 4.3.5 and 4.3.7]{PBC_ChenDeligne2017};
the latter proof reduces determinant one to the regular sequence
$(X-1,Y-1)$, and then to conjugation by a power of $c_0$.
Their determinant image is exactly the augmentation-one units,
which gives the stated quotient. Inner units act trivially on
this abelian quotient, so the $S_3$ action is substitution on $R$.

Finally our program expands the six standard Lyndon brackets of
length five and computes the integral kernel of
$(a,b)\mapsto[x,a]+[y,b]$ into degree six. It has the three
integral basis vectors recorded in its output. Their normalized
metabelian images $q=b-a$ are
$u^2(u+v),-2uv(u+v),-v^2(u+v)$.
The free coordinates of the kernel basis have coefficient one,
so this is its integral lattice, not merely a rational basis.
The resulting beta direction has the additional common factor
$uv$. These are all principal-block directions. A character
idempotent supported on nontrivial cubic characters has zero
image on that block, so the directions are not supplied by $e_0$.
\end{proof}

\begin{statement}{B.10}{Proved}{Finite beta cocycles are exactly the bounded root-valued coboundaries}
The subset $\mathcal T_1$ consists exactly of
\[
 \beta_g(X,Y)=\frac{g(X)g(Y)}{g(XY)},
\]
where $g$ is a finite-order unit of
$\widehat{\mathbb Z}[[\widehat{\mathbb Z}]]$,
$g(1)=1$, and $g(T^{-1})=g(T)^{-1}$. This $g$ is unique.
The order of $\beta_g$ is the order of $g$.
Thus Lead 7 is correct for nonzero $e_1$; for $e_1=0$ its claimed
non-inner involution is the identity and is not an involution.
\end{statement}
\begin{proof}
A finite-order one-variable unit is a root-valued function on the
prime-to-$p$ characters, constant in the pro-$p$ formal variable,
with Frobenius covariance. The formula makes $\beta_g$ symmetric
and invariant under replacing $(X,Y)$ by $(Y,(XY)^{-1})$.
Its boundary values are one. B.8 therefore puts it in $\mathcal T_0$,
and the cocycle identity follows by cancellation of the four
factors $g(X),g(Y),g(W),g(XYW)$. This proves the forward construction.
A bounded-order multiplicative function on the divisible torsion
group of prime-to-$p$ roots is trivial. Consequently a finite-order
$g$ with $\beta_g=1$ is one, which also proves uniqueness and the
assertion about orders once surjectivity is shown.

We prove that surjectivity, including the coefficient issue.
Fix $p$ and a common exponent $n$ for the values of $\beta$.
Let $D_p$ be the divisible group of all prime-to-$p$ roots of unity,
and let $C$ be the finite group of $n$-th roots in the maximal
unramified extension of $\mathbb Q_p$. The character values of
$\beta$ define a normalized symmetric two-cocycle $b:D_p^2\to C$,
with $b(a,a^{-1})=1$. It gives an abelian extension
$1\to C\to E\to D_p\to1$ with section $s(a)=(a,1)$.
First take the prime-to-$p$ part of $C$. Frobenius covariance gives
an automorphism $T$ of $E$ with $T(s(a))=s(a^p)$ and $T(c)=c^p$.
The difference between $T$ and the $p$-power homomorphism vanishes
on $C$ and hence is a homomorphism $D_p\to C$, which is zero.
Thus $T(e)=e^p$ for every $e\in E$.

For every prime $l\ne p$ and sufficiently large $k$, there is an
integer $m$ with $v_l(p^m-1)=k$. This follows by applying the
binomial expansion to a suitable power of $p$ which is $1$ modulo
$l$ (modulo $4$ when $l=2$). If $a$ has order $l^k$, covariance
now gives $s(a)^{p^m-1}=1$. Its $l$-primary part is therefore a
lift of $a$ of order dividing $l^k$. Choosing such lifts at higher
levels supplies every finite set of compatibility conditions for
a homomorphic section of $D_p[l^\infty]$; compactness of the finite
sets of possible lifts supplies a compatible family. Taking the
sum of these sections over the primary components splits the
extension. For $p=2$ the possible additional factor $C_2$ also
splits uniquely: $D_2$ is uniquely two-divisible, and the map
$a\mapsto 2\widetilde{a/2}$, in additive notation, is independent
of the lift and is a homomorphic section.

Every splitting is unique, since $\Hom(D_p,C)=0$. Uniqueness makes
it Frobenius-equivariant. Writing it relative to $s$ supplies the
function $g_p$ with $b=\delta g_p$; its values have order dividing
$n$. The equality $b(a,a^{-1})=1$ gives
$g_p(a^{-1})=g_p(a)^{-1}$. These functions over all primes define
the required finite-order unit, with their single common exponent.
This proves the exact classification. For the cubic example of
Lead 7, the saved character program verifies the cocycle identity
on all $216$ sixth-root triples, with no failures, and distinguishes
it from $\beta_0$.
\end{proof}

\subsection{Central lifting and the full lift of the cubic involution}

\begin{statement}{B.11}{Proved}{The general central criterion of Lead 26(a)}
Let $N$ be a closed characteristic subgroup of the free profinite
group $F$ of rank two. A positive peripheral automorphism of $F/N$,
with three prescribed conjugators, lifts peripherally to $F$ with
those conjugators if and only if it lifts in that manner to
$F/\overline{[N,F]}$. For $N=F''$ the central kernel is canonically
$H_2^{\mathrm{cts}}(M,\widehat{\mathbb Z})$.
\end{statement}
\begin{proof}
Necessity follows by quotienting. For sufficiency repeat the finite
argument of P.3 with an arbitrary finite source $F\twoheadrightarrow E$
and $A$ the image of $N$. Its scalar auxiliary source
$F\to D^\dagger$ is a central extension of $D=E/A$.
The image of $N$ in it is central, so this actual source map kills
$[N,F]$. Precomposition by the asserted automorphism of
$F/[N,F]$ is onto and retains the specified conjugator coordinates.
The subgroup-by-subgroup trace counts and Mobius inversion in P.3
therefore give a generating lift to $E$. Impose every finite source
condition and every finite conjugator coordinate together.
Compactness gives a full endomorphism generating every finite
quotient, hence an automorphism. This uses only normality of the
finite $A$, exactly as Lead 26(a) asserts.

The continuous homology five-term sequence for
$1\to F''\to F\to M\to1$ identifies its second homology with
$F''/\overline{[F'',F]}$, since $H_2(F,\widehat{\mathbb Z})=0$
and $H_1(F)\to H_1(M)$ is an isomorphism. All tensor products,
coinvariants and exterior powers below are completed, and all
commutator subgroups in this argument are closed.
\end{proof}

\begin{statement}{B.12}{Proved}{Every finite-order unitary metabelian peripheral class has a full lift}
Let $\alpha$ be a finite-order multiplier-one pure peripheral
outer automorphism of $M$, with beta unit $\beta$. If
\[
                    \beta(X,Y)\beta(X^{-1},Y^{-1})=1,
\]
then $\alpha$ lifts to the full profinite group $\mathcal P_1$,
with its prescribed normalized metabelian peripheral conjugators.
Every nontrivial such lift has infinite outer order.
In particular this holds for $\varphi_0$, every $\varphi_g$ of B.10,
and the entire product of cubic involutions in R28. Their central
obstruction in $H_2(M,\widehat{\mathbb Z})$ is zero, including at two.
\end{statement}
\begin{proof}
We compute the actual universal central kernel and keep its
integral torsion, since a rational or pro-prime lift alone would
not prove the assertion. Put $N=F''$, $E=F/[N,F]$, and
$C=N/[N,F]=H_2(M,\widehat{\mathbb Z})$.
For each prime $p$, let $C_p=H_2(M,\mathbb Z_p)$ and
$R_p=\mathbb Z_p[[A]]$. All modules, exterior products and
coinvariants in this proof are compact and completed.

In $1\to M'\to M\to A\to1$, the prime-to-$p$ part of $M'$
is homologically acyclic for $\mathbb Z_p$ coefficients, and its
pro-$p$ part is the torsion-free compact $\mathbb Z_p$-module $R_p$.
Thus $H_1(M',\mathbb Z_p)=R_p$ and
$H_2(M',\mathbb Z_p)=\widehat\bigwedge^2_{\mathbb Z_p}R_p$.
The regular module has $H_i(A,R_p)=0$ for $i>0$, as follows
by the completed induced-module resolution. Also
$H_3(A,\mathbb Z_p)=0$, since the pro-$p$ part of $A$ is
$\mathbb Z_p^2$ and its complementary part is acyclic. The
transgression $H_2(A,\mathbb Z_p)\to H_0(A,R_p)$ is an
isomorphism: both groups are $\mathbb Z_p$, and the class of
the two abelian generators maps to $[x,y]=c_0$, the module generator.
The degree-two Hochschild--Serre filtration therefore gives
\[
 C_p\simeq
 \bigl(\widehat\bigwedge^2_{\mathbb Z_p}R_p\bigr)_A.
\]
There are no omitted incoming differentials to this term: their
only possible sources are $H_2(A,R_p)=0$ and
$H_3(A,\mathbb Z_p)=0$.

Write $\iota(a)=a^{-1}$ on $A$ and extend this additively to $R_p$.
Diagonal coinvariants identify
$R_p\widehat\otimes R_p$ with $R_p$ by
$a\otimes b\mapsto a^{-1}b$. Passing to the exterior quotient
therefore gives the explicit compact module
\[
 C_p\simeq R_p/J_p,\qquad
 J_p=(1+\iota)R_p+\mathbb Z_p\cdot1.
 \tag{B.12.1}
\]
This is an additive submodule, not an ideal. To check the equality,
a diagonal tensor $a\otimes a$ maps to $1$, and the sum of
$a\otimes b$ and $b\otimes a$ maps to
$a^{-1}b+b^{-1}a$. Conversely these relations kill every diagonal
tensor by expansion. Finite sums of Dirac elements are dense,
and the displayed submodule is compact, hence closed, so the
calculation remains exact after completion.

We now prove that $C_p$ is torsion-free, including $p=2$.
For odd $p$, $J_p=R_p^{\iota}$ and the quotient is the
anti-invariant direct summand, so it is torsion-free. At two we
show the same equality without dividing by two globally. Interpret
$R_2$ as the compact module of $\mathbb Z_2$-valued finitely
additive measures on the clopen subsets of $A$. Inversion has the
single fixed point $1$, because $A=\widehat{\mathbb Z}^2$ is
torsion-free. Choose a decreasing cofinal sequence of invariant
clopen neighbourhoods $V_n$ of $1$. Each compact annulus
$V_n\setminus V_{n+1}$ has a free involution and admits a clopen
half $B_n$, with the other half $\iota B_n$. This follows by a
finite refinement of clopen neighbourhoods disjoint from their
translates.

For an invariant measure $\mu$, the values $\mu(V_n)$ all have
the same parity, since the measure of each annulus is
$2\mu(B_n)$. Choose $a\in\mathbb Z_2$ with that parity. Define
a measure $\nu$ to agree with $\mu$ on $B_n$, to vanish on
$\iota B_n$, and to have
$\nu(V_n)=(\mu(V_n)-a)/2$. These prescriptions are compatible:
the difference of two consecutive neighbourhood values is
$\mu(B_n)$. They define a measure on all clopen subsets, since
any clopen subset is a union of a subset of finitely many annuli
and possibly one sufficiently small $V_n$. They give
\[
                    \mu=(1+\iota)\nu+a\delta_1.
\]
Thus $R_2^{\iota}=J_2$. If $2r\in J_2$, invariance and
torsion-freeness of $R_2$ imply $r=\iota r$, hence $r\in J_2$.
This proves that multiplication by two is injective on $C_2$;
all odd integers are already units there. Therefore every $C_p$,
and consequently $C=\prod_p C_p$, has no nonzero finite-order
element.

Naturality gives the action of $\alpha$ on this central kernel.
It acts on $R_p=M'_p$ by multiplication by $\beta_p$. On the
diagonal tensor coinvariants its action is multiplication by
$\iota(\beta_p)\beta_p$, by the same formula
$a\otimes b\mapsto a^{-1}b$. Hence the unitary hypothesis makes
its action on every $C_p$, and on $C$, the identity.

Normalize $\alpha$ as $x\mapsto x$,
$y\mapsto f^{-1}yf$, where $f\in M'$, and choose any lift
$\widetilde f\in F'$ of $f$. This normalized representative has
the same order as its outer class. Indeed if its $n$th power is
$\operatorname{Inn}(c)$, fixing $x$ forces $c\in\langle x\rangle$
by B.7. Its conjugator at $y$ lies in $M'$. Comparing it with
conjugation by $c$, and using $C_M(y)=\langle y\rangle$, forces
the exponent of $x$ in $c$ to vanish in $M^{\mathrm{ab}}$.
Thus $c=1$. Equivalently, B.7's unique normalization is a group
section, and its restriction to multiplier one preserves order.
The endomorphism of $F$ given by
$x\mapsto x$, $y\mapsto\widetilde f^{-1}y\widetilde f$ induces
$\alpha$ on $M$, so preserves $N$ and $[N,F]$ and gives an
endomorphism $\Phi$ of $E$. On its quotient $M$ it is an
automorphism, and on its kernel $C$ it is identity by the preceding
natural homology calculation. The short exact sequence therefore
shows directly that $\Phi$ is an automorphism.

If $n$ is the order of $\alpha$, its conjugator norm
$f\alpha(f)\cdots\alpha^{n-1}(f)$ centralizes $y$ in $M$.
It belongs to $M'$, so is one by B.7. The corresponding norm of
$\widetilde f$ consequently lies in the central kernel of $E$.
Thus $\Phi^n$ fixes both $x$ and $y$, and is identity on $E$.
Let $h\in M'$ be the normalized conjugator of the third peripheral
$w=xy$, and choose any lift $\widetilde h\in E$. Initially there
may be an error
\[
                 \Phi(w)=\widetilde h^{-1}w\widetilde h\,a,
                    \qquad a\in C.
\]
The analogous metabelian norm of $h$ is one, so its lifted norm
is central. Iterating the display $n$ times, using $\Phi^n=1$
and identity on $C$, gives $a^n=1$. The proved torsion-freeness
of $C$ forces $a=1$. Therefore $\Phi$ is peripheral with the
prescribed conjugators. B.11 lifts it to the full free profinite
source.

For $\varphi_0$, simultaneous inversion preserves $e_0$ and
$\beta_0^2=1$, so the unitary condition holds. In the coordinates
of R7 the conjugators are
$r=(\beta_0-1)/((X-1)(Y-1))$ and
$s=X(Y-1)r/(XY-1)$, with $f=c_0^r,h=c_0^s$;
they satisfy $r+\beta_0r=s+\beta_0s=0$. For B.10,
$g(T^{-1})=g(T)^{-1}$ implies
$\beta_g^{\mathrm{inv}}=\beta_g^{-1}$ directly. The product of
any set of the cubic sign idempotents is likewise a unitary unit
of order at most two, including infinitely many coefficient primes.
This proves all the claimed zero obstructions. A nontrivial
metabelian outer image gives a nontrivial full outer class, which
has infinite order by B.13.
\end{proof}

\begin{statement}{B.13}{Known}{The torsion results of R25 and their exact consequence}
Every finite subgroup of $\mathcal P$ has order at most two;
every nonidentity element of such a subgroup has multiplier $-1$
at every prime. In particular $\mathcal P_1$ and the corresponding
multiplier-one subgroups of the two Grothendieck--Teichm\"uller
groups are torsion-free. Thus no nonidentity element of
$\mathcal T$ has a finite-order peripheral lift.
\end{statement}
\begin{proof}
These are precisely LM4 D.32 and D.34 in R25, whose full arguments
and supplied referee checks are distinguished from the unreviewed
rows of that return. The first bounds finite subgroups through
the two-adic multiplier; the second identifies the resulting sign
at every prime. Their conclusion applies to a finite-order lift
of a nontrivial metabelian multiplier-one class because such a
lift lies in $\mathcal P_1$. Nothing in this implication asserts
conjugacy of all involutions to $c$.
\end{proof}

\begin{statement}{B.14}{Proved}{The tripod centralizer and the four-cusp filling obstruction}
One has $C_{\mathcal P}(c)=\langle c\rangle$.
The four-cusp involution constructed in R22, LM4 D.29, cannot be
descended by filling one of its four punctures.
\end{statement}
\begin{proof}
A normalized pair $(\lambda,f)$ commuting with $c$ has
$f(x,y)=f(x^{-1},y^{-1})$. The fixed-group calculation in P.1
uses only the free profinite group, so gives $f=1$ without a
symmetry assumption. Evaluate $x,y$ at the two matrices $U,L$
of C.12. Their product has order six and trace one. Since
$\lambda$ is a profinite unit, $(UL)^\lambda$ also has trace one.
Peripherality therefore implies
$\operatorname{tr}(U^\lambda L^\lambda)=2-\lambda^2=1$ at every
prime. Hence $\lambda^2=1$, and the normalized power-substitution
automorphism has order at most two. B.13 identifies its multiplier
as the global $-1$ unless it is identity. This proves the
centralizer statement, with its torsion and multipliers.

For the second assertion the four actual peripheral generators in
LM4 D.29.7 are $(a,a^{-1}l,b,b^{-1}l^{-1})$, with product one.
The two factor maps in D.29.4--5 send their classes by the
permutation $(1\ 3)(2\ 4)$, as is verified directly from those
maps. Thus no individual filled peripheral normal subgroup is
preserved. Its normal closure differs from that of its partner
already in the abelianization of the free rank-three group.
A quotient obtained by filling one puncture consequently does
not carry the proposed descended action. This is a definite
failure of that descent, not evidence for an exotic tripod involution.
\end{proof}

\begin{statement}{B.15}{Proved}{The three-section generation criterion}
A negative peripheral involution of the full profinite tripod is
conjugate in $\mathcal P$ to complex conjugation if and only if
its three section classes admit involutions whose successive
products are the three marked inertia generators and which
generate the full extension $F\rtimes C_2$, with the chosen
orientation of those products. This is the criterion suggested
from cote 145, scan 23 and scans 28--29.
\end{statement}
\begin{proof}
For complex conjugation the three real-arc lifts are the three
free factors of $C_2\amalg C_2\amalg C_2$, and their successive
products generate its index-two free subgroup and the three
peripheral classes. Peripheral conjugacy transports this property.
Conversely the proposed three lifts define an epimorphism from
that free profinite product to the involution extension. It
preserves the quotient to $C_2$ and restricts to an epimorphism
between free profinite groups of rank two. After identifying their
marked free sources, the latter is an epimorphic endomorphism,
hence an isomorphism. The whole extension map is consequently an
isomorphism. Its restriction is peripheral by the stipulated
inertia products and intertwines the two order-two actions modulo
an inner automorphism. This is exactly conjugacy in $\mathcal P$.
The generating hypothesis cannot be replaced by generation in
every maximal pro-prime quotient.
\end{proof}

\begin{statement}{B.16}{Proved}{The order-75 residue already fails in one relative pro-two tower}
Let $q:F\twoheadrightarrow Q_{75}=\mathbb F_{25}\rtimes C_3$
be the quotient of R7, with
$q(x)=(0,1)$ and $q(y)=(1,1)$, and put
\[
 V=\ker q,\qquad N=\ker(V\longrightarrow V^{(2)}),
 \qquad G=F/N.
\]
There are no $f,h\in G'$ with the prescribed R7 residues for
which the substitution $\alpha(x)=x$, $\alpha(y)=f^{-1}yf$
satisfies all three equations
\[
 f\alpha(f)=1,\qquad
 \alpha(xy)=h^{-1}xyh,\qquad h\alpha(h)=1.
\]
Thus an obstruction occurs in a finite quotient of this specific
relative pro-two tower, rather than merely somewhere among the
finite quotients of $F$. This assertion does not specify its finite
depth; the explicit finite quotient is supplied separately in B.19.
\end{statement}
\begin{proof}
We first check the cohomological comparison that makes this
relative quotient usable. The kernel $N$ is two-perfect:
$H^1(N,\mathbb F_2)=0$. Indeed any continuous character on $N$
factors through its image in some finite quotient $D$ of $V$.
That image is $O^2(D)$, the smallest normal subgroup with
two-group quotient. This subgroup has no quotient of order two:
its subgroup generated by commutators and squares is characteristic,
and quotienting by it would still give a two-group quotient of $D$.
The minimality of $O^2(D)$ proves the assertion. The factorization
used here follows by choosing an open normal subgroup of $V$
whose intersection with $N$ is contained in the character's kernel.

For every open $U\subset G$, with inverse image $\widetilde U$
in $F$, inflation therefore gives
\[
 H^1(U,\mathbb F_2)\simeq H^1(\widetilde U,\mathbb F_2),
 \qquad \dim H^1(U,\mathbb F_2)=1+[G:U].
 \tag{B.16.1}
\]
In particular it retains every boundary class of the associated
finite covering of the tripod. The open subgroup $V/N$ is free
pro-two of rank $76$ and has odd index $75$. Hence
$\operatorname{cd}_2G=1$, $G$ has no element of order two, and
$\dim H^1(G,\mathbb F_2)=2$.

Suppose the three equations held. The first makes $\alpha^2=1$;
the residue $\overline\alpha(v,i)=(-v,i)$ makes its order exactly
two. In $E=G\rtimes\langle\alpha\rangle$ the three elements
$\alpha,f\alpha,h\alpha$ are involutions centralizing respectively
$x,y,xy$. The pro-two parts of these three peripheral groups are
nontrivial and are detected by the abelian characters of $G$.
Thus every one of the three associated fixed groups has nonzero
$H^1(-,\mathbb F_2)$.

The profinite Brown restriction theorem, in the precise form
\cite[Theorem 2 and \S1.1, equations (3)--(6)]{PBC_KZ2007},
applies to $E$: its open subgroup $G$ has two-cohomological
dimension one, and a finite two-subgroup injects into $E/G=C_2$.
It gives
\[
 \sum_{[T]}\bigl(1+\dim H^1(C_G(T),\mathbb F_2)\bigr)
       =\dim H^k(E,\mathbb F_2)\leq3\qquad(k>1),
 \tag{B.16.2}
\]
where $T$ runs over conjugacy classes of order-two subgroups.
The right-hand bound follows directly from Hochschild--Serre
and $H^i(G,\mathbb F_2)=0$ for $i>1$.
Since two of the three peripheral section classes would already
contribute at least four, those three classes coincide. Conjugacy
in $E$ is conjugacy by $G$ for sections of this split extension.
Consequently $C=C_G(\alpha)$ contains conjugates of the pro-two
parts of all three peripheral groups, and
$\dim H^1(C,\mathbb F_2)\leq2$.

The image of $C$ in $Q_{75}$ lies in the fixed subgroup of
$\overline\alpha$, which is $\{(0,i):i\in C_3\}$.
Let $U_0$ be its inverse image in $G$; it has index $25$.
By (B.16.1), $H^1(U_0,\mathbb F_2)$ has dimension $26$.
Restriction to $C$ has a nonzero kernel, so some open subgroup
$U\subset U_0$ of index two contains $C$.
The corresponding covering of the tripod has degree $50$.

Above each of the three base punctures choose the boundary
containing the stipulated conjugate of its pro-two inertia.
Each selected boundary has odd degree. There are additional
boundaries, since the sum of boundary degrees over each puncture
is the even number $50$. In the mod-two homology of a punctured
oriented surface, the boundary classes have only their total-sum
relation. Our three selected classes are therefore independent.
Their nonzero pro-two generators have the same mod-two classes,
since their boundary degrees are odd. By (B.16.1), the three dual
characters descend to $U$ and restrict independently to $C$.
This gives $\dim H^1(C,\mathbb F_2)\geq3$, a contradiction.

Finally the free pro-two kernel has a cofinal descending sequence
of characteristic open subgroups. The finite solution sets for the
three displayed equations, with their fixed $Q_{75}$ residues,
form an inverse system. If every one were nonempty, compactness
would give the excluded solution in $G$. Some finite level in this
one tower therefore fails. B.19 supplies an explicit finite quotient in this tower; its
certificate is separate from the cohomological argument (B.16.2).
\end{proof}

\begin{statement}{B.17}{Proved}{The exact central obstruction and the full metabelian image}
Put $R=\widehat{\mathbb Z}[[A]]$, write $v^*=v(X^{-1},Y^{-1})$,
and set $d=(X-1)(Y-1)(XY-1)$.
Let $\beta=1+dt\in R^\times$ be any multiplier-one pure
peripheral metabelian unit of B.7. In
\[
 C=F''/\overline{[F'',F]}\simeq R/R^*,
 \qquad R^*=\{v:v^*=v\},
\]
its central peripheral obstruction $a_\beta$ is uniquely determined by
\[
             2a_\beta=
       \left[-XY\,\frac{\beta\beta^*-1}{d}\right].
 \tag{B.17.1}
\]
Here the numerator is divisible by the regular element $d$;
the brackets denote an additive quotient, and division by two is
not imposed in the integral module. Equation (B.17.1) specifies
an existing integral element and is valid also at the coefficient
prime two. In particular
\[
 \operatorname{im}\bigl(\mathcal P_1\longrightarrow\Out(M)\bigr)
   =\{\beta\in(1+dR)\cap R^\times:\beta\beta^*=1\}.
 \tag{B.17.2}
\]
This is a statement about the full free profinite source and about
arbitrary units, with no finite-order assumption.
\end{statement}
\begin{proof}
The literature searches recorded with B.2 and in the work record
found the beta--gamma identities, but not this surjectivity
statement for the whole peripheral group. We give the central
calculation independently; no priority is claimed.

Let $E=F/\overline{[F'',F]}$ and use the central-kernel computation
of B.12. We fix its sign by specifying the bracket itself:
for $u,v\in R=M'$, arbitrary lifts to $E'$ satisfy
\[
             [\widetilde u,\widetilde v]=[u^*v]\in C.
 \tag{B.17.3}
\]
Indeed conjugating the commutator of the Dirac elements indexed
by $a,b\in A$ by $a^{-1}$ identifies it with the Dirac element
$a^{-1}b$ in the exterior-coinvariant calculation. Bilinearity,
density, and continuity give (B.17.3) for all $u,v$.

The normalized metabelian conjugator exponents are
\[
             r=(XY-1)t,\qquad s=X(Y-1)t.
 \tag{B.17.4}
\]
Choose any lifts $\widetilde f,\widetilde h\in E'$ of
$c_0^r,c_0^s$. The substitution $x\mapsto x$,
$y\mapsto\widetilde f^{-1}y\widetilde f$ induces an
endomorphism $\Phi$ of $E$. It is an automorphism: its map on
$M$ is the given metabelian automorphism, and its map on $C$
is multiplication by the invariant unit $\beta^*\beta$, by
B.12's natural tensor calculation. Its discrepancy is
\[
             \Phi(xy)=\widetilde h^{-1}xy\widetilde h\,a_\beta.
 \tag{B.17.5}
\]
Changing either lift by a central element leaves this equation
unchanged. Changing a metabelian frame by its peripheral
centralizer also leaves it unchanged: B.7 identifies that
centralizer as the corresponding procyclic subgroup, whose
powers lift to commuting powers of the same generator in $E$.
Thus no uncomputed centralizer-commutator quotient remains here.

We compute (B.17.5) using rational class-two coordinates, and
then descend the resulting integral identity. The group $E'$ is
nilpotent of class at most two, with torsion-free central kernel
$C$ and torsion-free quotient $R$. It is therefore torsion-free.
To construct the required rational coordinates explicitly, first
adjoin halves to its central kernel and define
$u\mathbin\oplus v=uv[u,v]^{-1/2}$. The class-two identities
make this an associative commutative operation and make the
commutator bilinear. Tensoring this torsion-free additive group
with $\mathbb Q$ gives a Lie algebra in which $E'$ embeds and
whose multiplication is $u*v=u+v+\tfrac12[u,v]$.
Its central sequence is the rationalization of
$0\to C\to E'\to R\to0$. Choose a merely linear splitting
$\sigma:R\otimes\mathbb Q\to\log E'_{\mathbb Q}$ with
$\sigma(1)=\log c_0$; continuity of this auxiliary splitting
is neither asserted nor needed.

Conjugation by $x,y$ fixes $C$ and takes the form
\[
 T_X\sigma(v)=\sigma(Xv)+k_X(v),\qquad
 T_Y\sigma(v)=\sigma(Yv)+k_Y(v).
\]
Since $[x,y]=c_0$, one has
$T_XT_Y=\operatorname{Ad}(c_0)T_YT_X$. Its central coordinate is
\[
 k_X((Y-1)v)-k_Y((X-1)v)=[XYv].
 \tag{B.17.6}
\]
The BCH products for the two sides of (B.17.5), with their common
factor $xy$ removed, now give
\[
\begin{split}
 a_\beta={}&k_X((Y-1)r-Ys)+k_Y(r-s)\\
           &-\tfrac12[r^*Yr]+\tfrac12[s^*XYs].
\end{split}
\]
Substitute (B.17.4) and (B.17.6). Because
$(Y-1)r-Ys=-(Y-1)t$ and $r-s=(X-1)t$, this becomes
\[
 a_\beta=
       \left[-XYt-\frac Y2 r^*r+\frac{XY}{2}s^*s\right]
       \quad\hbox{in }C\otimes\mathbb Q.
 \tag{B.17.7}
\]
Write $T=t^*t$. Straight expansion gives
\[
\begin{split}
 2a_\beta
   &=\left[-2XYt+(X^{-1}-X+2XY-2Y)T\right],\\
 -XY\frac{\beta\beta^*-1}{d}
   &=-XYt+(XY)^{-1}t^*+\frac d{XY}T.
\end{split}
\]
In $R/R^*$ the middle term of the second line equals
$-[XYt]$. The difference of the two right sides is
\[
       \bigl(XY+(XY)^{-1}-Y-Y^{-1}\bigr)T,
\]
which is invariant under $*$. This proves (B.17.1) rationally.
Both sides of (B.17.1) are integral, and B.12 proves that
$C\to C\otimes\mathbb Q$ is injective. Hence it proves exactly
the displayed integral identity, including at two.
The independent program \texttt{pbc\_obstruction.py} checks the
Laurent identities and fifteen exact peripheral discrepancies in
independently constructed Heisenberg coordinates.

Finally put $P=-XY(\beta\beta^*-1)/d\in R$.
The identities $d^*=-d/(XY)^2$ and
$(\beta\beta^*)^*=\beta\beta^*$ give $P^*=-P$.
An anti-invariant element lying in the invariant submodule is
zero, since $R$ is torsion-free. Consequently
$[P]=0$ if and only if $P=0$, equivalently
$\beta\beta^*=1$. By torsion-freeness of $C$ and (B.17.1),
this is also equivalent to $a_\beta=0$.
B.11 now gives the full peripheral lift in precisely these cases.
Conversely any full peripheral lift descends to a peripheral
automorphism of $E$, so has zero discrepancy and satisfies the
same condition. This proves both directions of (B.17.2).
\end{proof}

\begin{statement}{B.18}{Proved}{All members of the two R28 families and their lift obstruction}
The families of LM6 A.10 and A.49 exist as stated in the full
metabelian group. For A.49 write
\[
       u(\mathbf t)=1+\sum_p(t_p-1)e_p,
       \qquad \mathbf t=(t_p)_p\in\prod_p\mathbb Z_p^\times,
\]
with the cubic idempotents at $p\ne3$ and the six regular
fourth-root character pairs at $p=3$ specified below.
Its exact central obstruction has coefficient-$p$ component
\[
       a_{\mathbf t,p}
         =\left[-\frac{XY(t_p^2-1)}{2d}e_p\right]
           \in R_p/R_p^*.
 \tag{B.18.1}
\]
Therefore this member lifts to $\mathcal P_1$ if and only if
$t_p\in\{1,-1\}$ for every prime $p$. This includes all of A.10,
the coefficient-two sign, and the coefficient-three sign of A.49.
Every nonidentity lift of these signs has infinite outer order.

The R7 class $\varphi_0$ is the member $t_5=-1$, $t_p=1$ for
$p\ne5$. The cubic example $\varphi_g$ in Lead 7 is a different
arithmetic-centralizing involution and is outside the A.49 family.
Every finite $\varphi_g$ of B.10 nevertheless has zero central
obstruction and a full peripheral lift.
\end{statement}
\begin{proof}
For $p\ne3$, average over the kernel of
$A^{(p')}\to C_3$, $X,Y\mapsto\zeta_3$, and select the two
nontrivial characters. The compatible finite averages have
prime-to-$p$ denominators and give a genuine continuous
idempotent $e_p\in R_p$. Its character support is
$(\zeta_3,\zeta_3)$ and $(\zeta_3^2,\zeta_3^2)$, allowing
arbitrary pro-$p$ variables. For $p=3$ average over the kernel
of $A^{(3')}\to(\mathbb Z/4)^2$ and select the six characters
\[
    (a,b)\in(\mathbb Z/4)^2,
             \qquad a\ne0,\quad b\ne0,\quad a+b\ne0.
\]
Their Fourier sum is defined over $\mathbb Z_3$, since conjugate
characters occur together and $16$ is invertible there.
In all cases $d$ is invertible on $e_pR_p$.
The supports are invariant under the peripheral $S_3$, under
simultaneous inversion, and under every scalar unit of
$\widehat{\mathbb Z}$. Thus the displayed $u(\mathbf t)$ is an
element of B.7's group. Its actual normalized conjugators are
\[
 r(\mathbf t)=\frac{u(\mathbf t)-1}{(X-1)(Y-1)},
 \qquad s(\mathbf t)=\frac{X(Y-1)r(\mathbf t)}{XY-1}.
\]
The product formula of B.7 gives
$\varphi_{\mathbf t}\varphi_{\mathbf v}=\varphi_{\mathbf t\mathbf v}$.
The map is continuous, because any finite target uses only
finitely many coefficient primes and residue digits. It is
injective by the uniqueness of B.7's normalization, and hence
has closed image by compactness. This rederives both families,
including their coefficient-two and coefficient-three components.

For clarity, arithmetic commutation here is outer commutation.
An arithmetic representative acts on $A$ by the scalar
$\chi(\sigma)$. Conjugation on the abelian IA outer group is
therefore the natural substitution
$(X,Y)\mapsto(X^{\chi(\sigma)},Y^{\chi(\sigma)})$; the
arithmetic representative's own IA unit cancels. This also follows
from \cite[Proposition 4.3.7]{PBC_ChenDeligne2017}.
The invariant supports just checked prove commutation with every
$G_K$. Every maximal pro-prime quotient kills the nonprincipal
prime-to-$p$ support in its coefficient factor, and kills all
other coefficient factors. Hence these automorphisms are trivial
on every such quotient, exactly as asserted in R28.

Since $e_p^*=e_p$, one has
$u(\mathbf t)u(\mathbf t)^*-1=(t_p^2-1)e_p$ in the
coefficient-$p$ factor. Formula (B.17.1) gives (B.18.1).
Its representative is integral even at two, because
$(t_2^2-1)/2\in\mathbb Z_2$ for every $t_2\in\mathbb Z_2^\times$.
It is anti-invariant and is nonzero precisely when $t_p^2\ne1$.
The ring $\mathbb Z_p$ is a domain, so this condition vanishes
precisely at $t_p=\pm1$. B.17 proves the claimed complete lift
classification, and B.13 gives the infinite order of nontrivial
full lifts.

For the cubic $g$ of Lead 7, evaluate at the character pair
$(X,Y)=(\zeta_3,-1)$ in the coefficient-five factor. The value of
$g(X)g(Y)/g(XY)$ is $-1$, whereas every A.49 member has value
$1$ there. Thus this $\varphi_g$ is not a member of that family,
even though it agrees with $\varphi_0$ on the cubic diagonal
support. Its one-variable support is invariant under every scalar
unit, so the same arithmetic calculation puts it in every
metabelian arithmetic centralizer. Finally
$g(T^{-1})=g(T)^{-1}$ gives
$\beta_g^*=\beta_g^{-1}$ for all of B.10, so (B.17.1) is zero.
For a general finite $g$, arithmetic commutation over a specified
$K$ is the exact additional condition
$g(T^{\chi(\sigma)})=g(T)$ for every $\sigma\in G_K$:
necessity follows from B.10's uniqueness of the finite coboundary
primitive. It is not implied by inversion covariance alone.
\end{proof}

\begin{statement}{B.19}{Proved}{An explicit finite obstruction of order 614400}
There is an explicitly specified quotient
\[
                F\twoheadrightarrow G,
                \qquad |G|=75\cdot2^{13}=614400,
\]
above the R7 quotient $Q_{75}$, in which the prescribed
involutive lifting system has no solution. Every hypothetical full
solution with metabelian residue $\varphi_0$ necessarily preserves
this particular quotient; this descent is proved below, rather
than assumed from its order or its construction.
The quotient has an extraspecial two-group kernel over $Q_{75}$,
and both chosen generators have order three. Thus the obstruction
belongs to the relative pro-two tower of B.16.
\end{statement}
\begin{proof}
\emph{The group and its generators.}
Let
\[
 k=\mathbb F_2[\eta]/(\eta^4+\eta+1),\quad
 k_0=\mathbb F_4\subset k,
 \quad \zeta=\eta^3,
 \quad W=k^3.
\]
Then $\zeta$ has order five and $\eta+\eta^4=1$.
For $u,v\in W$ put
\[
\begin{split}
 F_0(u,v)&=\operatorname{Tr}_{k/\mathbb F_2}
                       \left(\eta\sum_{j=0}^2u_jv_j^4\right),\\
 Q_0(u)&=F_0(u,u)
        =\operatorname{Tr}_{k_0/\mathbb F_2}
                       \left(\sum_{j=0}^2u_ju_j^4\right),\\
 B_0(u,v)&=F_0(u,v)+F_0(v,u)
        =\operatorname{Tr}_{k/\mathbb F_2}
                       \left(\sum_{j=0}^2u_jv_j^4\right).
\end{split}
\]
The bilinear form $B_0$ is alternating and nondegenerate, by
nondegeneracy of the finite-field trace pairing. Define
\[
 P=W\times\mathbb F_2,
 \qquad (u,c)(v,d)=(u+v,c+d+F_0(u,v)).
 \tag{B.19.1}
\]
Bilinearity proves associativity. The square is
$(u,c)^2=(0,Q_0(u))$, and the commutator is
$(0,B_0(u,v))$. Hence $P$ has order $2^{13}$ and its center
and derived subgroup both have order two.

Write $Q_{75}=V_5\rtimes C_3$, where $V_5=\mathbb F_5^2$ and
the generator of $C_3$ acts by
\[
       T(a,b)=(-b,a-b).
\]
This identifies $a+b\omega\in\mathbb F_{25}$ with $(a,b)$,
where $\omega^2+\omega+1=0$. Put
$\ell_j(v)=(T^{-j}v)_1$, with indices modulo three, and define
\[
       (\rho(v,i)u)_j
            =\zeta^{\ell_j(v)}u_{j-i}.
 \tag{B.19.2}
\]
The identity $\ell_{j-1}(v)=\ell_j(Tv)$ verifies the semidirect
product law. These maps preserve $F_0$ because every diagonal
factor has fifth power one and the other operation permutes the
coordinates. They act on $P$ with its central coordinate fixed.
Let $G=P\rtimes_\rho Q_{75}$ and take
\[
 x_G=((0,0),((0,0),1)),\qquad
 y_G=(((\zeta,0,1),0),((1,0),1)).
 \tag{B.19.3}
\]
Both have order three. Direct multiplication gives, for $w=x_Gy_G$,
\[
 w^3=(v,0),\qquad
 v=(1+\zeta,1+\zeta,1+\zeta^4),\qquad Q_0(v)=1.
 \tag{B.19.4}
\]
For the last equality the first two norms cancel, while the last
is $\zeta+\zeta^4\in k_0$ of trace one. In particular $w$ has
order twelve.

The $\mathbb F_2[Q_{75}]$-module $W$ is irreducible.
On restriction to $V_5$, its three coordinate characters have
distinct Frobenius orbits, each of length four; $C_3$ cycles those
three simple summands. Thus a nonzero $Q_{75}$-submodule is all
of $W$. The images of $x_G,y_G$ generate $Q_{75}$, and their
subgroup contains $w^3$ with nonzero image $v$ in $W$.
Conjugating it and taking products gives all of $W$ in the image
of their intersection with $P$. Nondegeneracy of $B_0$ then puts
the central involution in that subgroup as well. Consequently
$x_G,y_G$ generate $G$, proving the asserted epimorphism from $F$.
Also $G'=P\rtimes V_5$ and $G/G''=Q_{75}$.

\emph{All possible induced automorphisms.}
The R7 automorphism of $Q_{75}$ is
\[
                  \overline\alpha(v,i)=(-v,i).
\]
The same character calculation shows
$\operatorname{End}_{\mathbb F_2[Q_{75}]}(W)=k$:
each of the three distinct $V_5$-summands has endomorphism field
$k$, and cyclic permutation makes their three scalars equal.
Entrywise fourth power intertwines $\rho$ with
$\rho\overline\alpha$. Hence every invertible intertwiner is
$L_a(u)=a u^4$ for some $a\in k^\times$.
It preserves $Q_0$ precisely when $aa^4=1$, namely when
$a\in\mu_5=\langle\zeta\rangle$.

For each such $a$ the following is an actual involution of $G$:
\[
 \Psi_a(u,c;q)=(a u^4,c+Q_0(u);\overline\alpha(q)).
 \tag{B.19.5}
\]
Indeed
$F_0(L_au,L_av)+F_0(u,v)=B_0(u,v)$, which is the polar form
of $Q_0$; this verifies the multiplication law, and applying the
formula twice gives identity.
Every automorphism of $G$ over $\overline\alpha$ is
\[
                        \operatorname{Inn}(k)\Psi_a,
                  \qquad k\in W,\quad a\in\mu_5.
 \tag{B.19.6}
\]
To prove exhaustiveness, remove (B.19.5). An automorphism inducing
identity on $W$ and on $Q_{75}$ is a central automorphism on $P$.
It is inner by a unique $k\in W$, because $B_0$ is nondegenerate.
After removing that inner map it fixes $P$ pointwise. Its remaining
change on $Q_{75}$ is a homomorphism to $Z(P)=C_2$, and is zero
because $|Q_{75}|$ is odd. This proves (B.19.6).
It fixes $x_G$ and has order dividing two exactly when
\[
       k=(b,b,b),\qquad a b^4=b.
 \tag{B.19.7}
\]
For each of the five $a$ there are four such $b$, including zero.
Thus there are exactly twenty candidates.

\emph{Why every hypothetical full solution descends here.}
This point is essential: arbitrary quotients of order $614400$
need not be preserved by a putative lift. A full solution of
(A--I) first descends to
\[
       T_0=C_3\amalg C_3
          =F/\langle\!\langle x^3,y^3\rangle\!\rangle,
\]
because the two cubes are taken to conjugates of themselves.
Its order is two and its $Q_{75}$ residue is
$\overline\alpha$. Let $U=\ker(T_0\to Q_{75})$ and
$D=U/\overline{U^2[U,U]}$.
The Bass--Serre covering graph has $75$ edges and two sets of
$25$ vertices, so $U$ is free profinite of rank $26$.
Its mod-two graph sequence gives, in the semisimple
$\mathbb F_2[Q_{75}]$-module category,
\[
 [D]=[\mathbb F_2[Q_{75}]]
     -2[\mathbb F_2[Q_{75}/C_3]]+[\mathbb F_2].
 \tag{B.19.8}
\]
For precision, the simple modules are the trivial line, the
two-dimensional nontrivial $C_3$ module $L$, and two modules
$W,W'$ of dimension twelve. The $24$ nonzero characters of
$V_5$ have two orbits under scalar Frobenius $2$ and the
order-three matrix $T$, each of length twelve. Each of $W,W'$
has endomorphism field $k$. Thus
\[
\begin{split}
 \mathbb F_2[Q_{75}]&\simeq1\oplus L\oplus W^{\oplus3}
                                         \oplus(W')^{\oplus3},\\
 \mathbb F_2[Q_{75}/C_3]&\simeq1\oplus W\oplus W',\qquad
 D\simeq L\oplus W\oplus W'.
\end{split}
\]
Every summand of $D$ therefore occurs once. The map
$T_0\to G/Z(P)=W\rtimes Q_{75}$ selects its unique $W$
quotient. The scalar $-1$ on $V_5$ is Frobenius $4$ on
its character values, so $\overline\alpha$ preserves each of
$W,W'$, rather than interchanging them. Consequently the kernel
of this particular quotient of $T_0$ is preserved by the
hypothetical involution. It induces an involution
$\gamma$ of $W\rtimes Q_{75}$.

Its map on $W$ is $L_a$ for some $a\in k^\times$.
Since $\gamma^2=1$, one has $aa^4=1$.
The remaining change of a complement $Q_{75}$ is a cocycle with
values in $W$. Averaging over the odd group $Q_{75}$ makes this
a coboundary. Hence $\gamma$ lifts to an automorphism of $G$
of the form (B.19.6). Such a lift is unique:
an automorphism over the central quotient differs by a continuous
homomorphism $W\rtimes Q_{75}\to C_2$, and this group has none,
since $W_{Q_{75}}=0$ and $Q_{75}$ has odd order. Uniqueness also
makes the lifted automorphism an involution.

Finally every element of order three in $G/Z(P)$ has a unique
lift of order three in $G$: its two lifts differ by the central
involution, and exactly one of their cubes is one.
The two homomorphisms from $T_0$ obtained by applying the
hypothetical source involution or this lifted target involution
have the same central quotient and send both factor generators
to elements of order three. They agree on those generators and
hence agree on $T_0$. This proves forced descent to $G$ itself.
It also transports both specified conjugators and both norm
equations, not just the generator images.

\emph{The finite norm contradiction.}
The required third frame has $Q_{75}$ residue
$h_0=((4,3),0)\in V_5$, as in R7. Suppose one of the induced
involutions $\alpha$ had such an $h$ with
$\alpha(w)=h^{-1}wh$ and $h\alpha(h)=1$.
Then $h\alpha$ is an involution in the split extension and
centralizes $w$, hence $w^3$. Its action on $P$ has the form
\[
        (u,c)\longmapsto
          (Su,c+Q_0(u)+B_0(l,Su)),
 \qquad
        (Su)_j=a\zeta^{\ell_j(h_0)}u_j^4
 \tag{B.19.9}
\]
for some $l\in W$. Each diagonal coefficient has norm one.
The fixed space $W^S$ is totally isotropic for $B_0$:
for two fixed vectors, each $u_jv_j^4$ is fixed by fourth
power and therefore belongs to $k_0$, whose trace from $k$
to $\mathbb F_2$ is zero.

Squaring (B.19.9) and using nondegeneracy of $B_0$ forces
$Sl=l$. If (B.19.9) fixes an element with vector coordinate
$v$, it then has $Sv=v$ and
$Q_0(v)+B_0(l,v)=0$. Total isotropy gives $Q_0(v)=0$.
But (B.19.4) supplies a fixed element $w^3$ with $Q_0(v)=1$.
This contradiction proves that the third peripheral and norm
conditions cannot hold in $G$.

The complete independent finite calculation is
\texttt{pbc\_extraspecial\_c3c3.py}. It checks the field,
bilinear forms, all $75$ representation matrices and their
multiplication, the semilinear maps, the generator orders,
and generation. For the pair (B.19.3), all twenty candidates
have eight admissible first frames, counting the central bit,
and none has an admissible third frame. Every one of the
$4096$ possible vector coordinates is tested for each frame;
the central bit cancels from both conjugacy and norm tests.
The file \texttt{pbc\_extraspecial\_norms.out} gives the shorter
norm certificate: only $a=\zeta^2$ admits any third conjugator,
and its sixteen possibilities have nonzero vector norm.
The symbolic contradiction above explains these exhaustive data.

For the characteristic finite-system formulation of R6 one can
use the following completely specified finite envelope, without
enumerating its much larger order:
\[
 C=F/\bigcap_{q\in\operatorname{Hom}_{\rm cont}(F,G)}\ker q.
 \tag{B.19.10}
\]
There are at most $|G|^2$ homomorphisms in this intersection.
Thus $C$ is finite, its kernel is fully invariant, and the
displayed generating pair gives a surjection $C\to G$.
An (A--I) solution in $C$ would descend to
$C_0=C/\langle\!\langle x^3,y^3\rangle\!\rangle$, which
is a quotient of $T_0$ and still maps onto $G$.
The mod-two abelianization of its $Q_{75}$-kernel is a quotient
of $D$ in (B.19.8). It therefore still has a unique $W$ quotient,
and the same multiplicity-one argument forces descent to
$W\rtimes Q_{75}$. The unique lift to $G$ and the unique
order-three lifts of the two generators then force descent to
$G$ with both frames and their norms. The contradiction already
proved applies. Consequently (A--I) is unsolvable in this
explicitly specified characteristic finite quotient as well.
\end{proof}

\subsection{Independent audit of the arithmetic derived-level assertions}

\begin{statement}{B.20}{Proved}{No arithmetic-centralizing IA image survives from length three to two}
Let $K$ be any number field and $S_j=F/F^{(j)}$, with the
arithmetic outer action of the tripod over $K$. If an element of
$C_{\Out(S_3)}(G_K)$ acts trivially on $F^{\mathrm{ab}}$, its
image in $\Out(M)$ is trivial. No peripheral hypothesis is needed.
This verifies LM6 A.56, including its application to every
nonidentity member of R28, while leaving the kernel of
$\Out(S_3)\to\Out(M)$ uncomputed.
\end{statement}
\begin{proof}
We record all comparisons needed for the supplied fibre-isolation
argument, rather than treating the previous return as authority.
First $S_j$ is center-free for $j\geq2$. The case $M$ follows
from its regular Alexander module. Inductively a central element
of $S_j$ lies in $F^{(j-1)}/F^{(j)}$. For each open normal
$H\supset F^{(j-1)}$, its class in the torsion-free
$H^{\mathrm{ab}}$ is $F/H$-invariant and maps to zero in
$F^{\mathrm{ab}}$. Transfer therefore kills it by $[F:H]$, so
the class is zero. The intersection of these $H'$ is exactly
$F^{(j)}$: separate an element in a finite quotient of
$F/F^{(j)}$ and take the inverse image of the abelian image of
$F^{(j-1)}$. This proves the induction.
For a center-free geometric group $S_j$, the arithmetic extension
is canonically
\[
 E_j\simeq\Aut(S_j)\mathbin\times_{\Out(S_j)}G_K,
\]
by conjugation and the projection to $G_K$. An outer centralizer
class therefore gives an actual automorphism $\beta$ of $E_j$
over $G_K$, after a geometric representative is chosen.

Write $R=\widehat{\mathbb Z}[[A]]$, with augmentation ideal $I$.
Chen--Deligne's IA description represents the metabelian image
by $u\in R^\times$ of augmentation one, modulo the monomials
$A$ \cite[Proposition 4.3.7]{PBC_ChenDeligne2017}.
The homomorphism
\[
        \partial:1+I\longrightarrow I/\overline{I^2}\simeq A,
           \qquad u\longmapsto u-1
\]
restricts to identity on $A$. The isomorphism follows both ways
by the square-zero ring $\widehat{\mathbb Z}\oplus A$.
Arithmetic commutation says $u/\sigma_R(u)\in A$, where
$\sigma_R$ is scalar cyclotomic substitution. Removing the
monomial with derivative $\partial u$ gives a representative
$u_0$ with derivative zero. Then $u_0/\sigma_R(u_0)$ is a
monomial with derivative zero, and is one. Hence $u_0$ is exactly
arithmetic-invariant.

Choose the geometric representative on $S_3$ whose metabelian
determinant is $u_0$. For an arithmetic lift $s(\sigma)$ write
$\beta(s(\sigma))=d_\sigma s(\sigma)$. Its commutator action
on $M'$ is multiplication by $u_0/\sigma_R(u_0)=1$.
The inner correction $d_\sigma$ consequently has trivial
abelian monomial, so lies in $M'$. Thus
\[
                       \beta|_{E_1}=1
 \tag{B.20.1}
\]
as an actual arithmetic automorphism, including the chosen
tangential section and all constant Kummer twists.

If $u_0$ is nontrivial modulo $A$, the closed cyclic subgroups
$\widehat{\mathbb Z}c_0$ and
$\widehat{\mathbb Z}u_0c_0$ in $M'$ are not conjugate.
Otherwise $u_0=ta$ for a scalar unit $t$ and a monomial $a$;
augmentation forces $t=1$. Compactness detects this failure in
a finite quotient. Refining to a characteristic quotient is
possible because $M$ has only finitely many open subgroups of
each index. We obtain a finite metabelian group $Q$, an induced
automorphism $\gamma$, and nonconjugate equal-order cyclic
subgroups $J=\langle g\rangle$ and $J'=\gamma(J)$ such that
\[
                      |(Q/J)^g|>0,\qquad |(Q/J')^g|=0.
 \tag{B.20.2}
\]
The zero follows because containment of a generator would make
the equal-order cyclic subgroups conjugate.

After a finite field extension the quotient is a constant
$Q$-torsor $Y\to\mathbb P^1-\{0,1,\infty\}$, and its
arithmetic quotient is $Q\times G_K$. The induced automorphism
has the form $(q,\sigma)\mapsto(\gamma(q)z(\sigma),\sigma)$,
where $z:G_K\to Z(Q)$ is a continuous homomorphism. A further
finite extension kills $z$; then take a Galois closure over
$\mathbb Q$. These operations retain (B.20.1).
Hilbert irreducibility supplies $a\in K-\{0,1\}$ whose
$Y$-fibre is a field with Galois group $Q$
\cite[Proposition 3.3.1 and Theorem 3.4.1]{PBC_SerreTopics}.
In its normal closure over $\mathbb Q$, extend $g$ to an
automorphism fixing $K$. Chebotarev gives a rational prime $q$
and a degree-one place of $K$ at which the fibre Frobenius is
$g$ \cite[Chapter V, Theorem 3.23]{PBC_MilneCFT}.
Exclude primes of bad marked reduction, divisors of $|Q|$, and
primes at which $a$ or $1-a$ is not a unit.

Set $n=q-1$ and enlarge constants to $K(\mu_n)$; the chosen
place still has residue field $\mathbb F_q$. The abelian cover
\[
        B_a:\quad z_1^n=t/a,\qquad
                     z_2^n=(1-t)/(1-a)
\]
is fixed by $\beta$ including its constants, by (B.20.1).
Its only rational fibre is over the reduction of $a$, and that
fibre has $n^2$ points. Therefore
\[
 \#((Y/J)\times B_a)(\mathbb F_q)=n^2|(Q/J)^g|>0,
 \quad
 \#((Y/J')\times B_a)(\mathbb F_q)=0.
 \tag{B.20.3}
\]
Both products have metabelian geometric monodromy. Every component
subgroup $H\subset F$ contains $F''$, hence $F'''\subset H'$.
The third derived arithmetic extension thus retains the full
abelianization of every component, with the permutation of the
components and their conjugation maps. The automorphism $\beta$
identifies their $H^0$ and $H^1$ arithmetic representations.
All the marked products have good reduction at the chosen place:
locally at a boundary point their normalized products have tame
Kummer equations with ramification index prime to $q$.
Proper-smooth base change, applied to compactifications and their
finite boundary sets, transfers these representations to the
special fibres. Poincar\'e duality identifies $H_c^1$ as well;
$H_c^2$ is their component module with the same Tate twist.
The affine trace formula then forces equality in (B.20.3), a
contradiction \cite[Theorems 20.4, 24.1, 29.4]{PBC_MilneLEC}.
This proves the assertion, with no connectedness assumption on
the two auxiliary products and no discarded arithmetic twist.
\end{proof}

\begin{statement}{B.21}{Proved}{The peripheral arithmetic projection is exactly geometric}
For every number field $K$, a pure peripheral outer automorphism
of $M$ commuting with $G_K$ has multiplier one. Moreover
\[
 \operatorname{im}\left(
 C_{\Out_{\partial}(S_j)}(G_K)\longrightarrow\Out(M)\right)
                  =\mathfrak S_3\qquad(j\geq3),
\]
where $\Out_{\partial}$ allows permutation of the three peripheral
classes. This verifies LM6 A.57 and A.58. It does not assert that
the same-level centralizer has trivial kernel.
\end{statement}
\begin{proof}
Remove a geometric cusp permutation. A pure class acts on $A$
by a scalar $\rho$. The center-free construction in B.20 gives
an arithmetic automorphism of $E_2$.
The geometric normalizers of $\langle x\rangle$ and
$\langle y\rangle$ in $M$ are those same groups: normalizing
forces exponent one by abelianization, then B.7 gives the
centralizer. The unit tangential sections at $0$ towards $1$
and at $1$ towards $0$ have the same image in
$E_1=A\rtimes G_K$. Indeed the Kummer covers
$t^{1/n},(1-t)^{1/n}$ have leading coefficient or value one,
with the compatible rational root labels one.

Write the induced map on $E_1$ as
$(v,g)\mapsto(\rho v+c(g),g)$.
Preservation of the two arithmetic normalizers forces the second
coordinate of $c$ to be the coboundary supplied by the first
cusp conjugator, and the first coordinate to be that supplied
by the second. Explicitly, if their abelian coordinates are
$a_0,a_1$, then
$c=(1-\chi)((a_1)_1,(a_0)_2)$.
An inner adjustment therefore makes the map exactly
$(v,g)\mapsto(\rho v,g)$.

Choose an odd rational prime $q$ splitting in a Galois closure
of $K$, put $n=q-1$, and pass to $K(\mu_n)$ at a degree-one
place. For every $a\in\{2,\ldots,q-1\}$ the cover $B_a$ used
in B.20 is taken to
\[
       w_1^n=t/a^r,\qquad
       w_2^n=(1-t)/(1-a)^r,
            \quad r=\rho\pmod n.
\]
This follows directly from the arithmetic Kummer quotient
$(v,g)\mapsto v-(\kappa_n(a)(g),\kappa_n(1-a)(g))$.
The geometric monodromy is abelian, so $E_2$ already retains
the full $H^1$ of every component: $H\supset F'$ implies
$F''\subset H'$. The good-reduction and trace comparison in
B.20 therefore makes the two point counts equal. The first is
$n^2$; the second is $n^2$ exactly when
\[
                         a^r+(1-a)^r=1\quad\hbox{in }\mathbb F_q.
\]
This identity for every $a$ makes the multiplicative bijection
$u\mapsto u^r$ additive: substitute $a=u/(u+v)$, and handle
$u+v=0$ using oddness of $r$. It is an automorphism of the
prime field, hence identity. Thus $\rho\equiv1\pmod{q-1}$.
Choosing $q$ also split in $\mathbb Q(\mu_m)$ for arbitrary
$m$ gives $\rho=1$. This proves the first assertion.
For the second, remove the geometric permutation and apply
B.20 to the remaining IA class. The reverse inclusion is the
geometric action, defined over $\mathbb Q$ at every derived level.
\end{proof}

\begin{statement}{B.22}{Proved}{One further derived layer detects all element conjugacy classes}
For $m\geq2$, the projection to $\Out(S_m)$ of a peripheral
arithmetic centralizer class at level $m+1$ is class-preserving,
after its unique geometric cusp permutation has been removed.
Every continuous class-preserving automorphism of $S_2$ or
$S_3$ is inner. Consequently
\[
 \operatorname{im}\bigl(C_{\Out_{\partial}(S_j)}(G_K)
                 \longrightarrow\Out(S_3)\bigr)
                  =\mathfrak S_3\qquad(j\geq4).
\]
These are the precise projected conclusions of LM6 A.66 and
A.67. They leave the kernel at the source level undetermined.
\end{statement}
\begin{proof}
By B.21, remove the unique cusp permutation and choose a
geometric representative acting literally as identity on $M$.
Its arithmetic correction has central image in $M$, hence zero.
Thus the arithmetic automorphism of $E_{m+1}$ acts literally
as identity on $E_2$ and $E_1$.
Take any characteristic finite quotient $Q$ of $S_m$. As in
B.20, make its deck action constant and kill the finite central
arithmetic twist by a field extension. For any $g\in Q$,
Hilbert irreducibility and degree-one Chebotarev then give a
good rational prime $q$ and a base fibre at $a$ with Frobenius
$g$. The cover $B_a$ with exponent $q-1$ isolates that fibre.
For $Z=Y\times B_a$, every geometric component subgroup
$H$ contains $F^{(m)}$, so $F^{(m+1)}\subset H'$.
The arithmetic automorphism at level $m+1$ therefore identifies
the complete component and first-homology diagrams of $Z$,
intertwining deck transformation $h$ with $\alpha_Q(h)$.
The base-change, duality, and trace argument of B.20 applies
also to deck-twisted Frobenius. Such a trace counts fixed points
of $h\operatorname{Frob}$: equivalently twist the finite-field
descent datum by the constant finite-order deck transformation.
On the isolated regular $Q$-fibre, deck action is left
multiplication and Frobenius is right multiplication by $g$.
Consequently the number of fixed points is
\[
 (q-1)^2
 \begin{cases}|C_Q(g)|,&h\text{ conjugate to }g^{-1},\\
 0,&\text{otherwise}.
 \end{cases}
\]
Comparing $h=g^{-1}$ and $\alpha_Q(h)$ proves that
$\alpha_Q(g)$ is conjugate to $g$. For each fixed element of
$S_m$, its nonempty closed sets of conjugators in all
characteristic finite quotients have the finite intersection
property. Compactness supplies a conjugator in $S_m$ itself.
This establishes the class-preserving assertion; no uniform
conjugator has yet been inferred.

Here are the completed-module details that supply that stronger
conclusion at levels two and three. Put
$\Lambda_P=\widehat{\mathbb Z}[[P]]$ for a profinite group $P$.
If $u\in P$ has infinite $p$-primary order, multiplication by
$u-1$, on either side, is injective on $\mathbb F_p[[P]]$.
Indeed invariance makes coefficients constant on the $u$-orbits
in each finite quotient. Refine the quotient so every orbit
length is multiplied by $p$; projection of each invariant orbit
sum is then zero. Such refinements exist from the induced
topology of the closed cyclic subgroup. Successive division by
$p$ proves the same fact over $\mathbb Z_p$.

If commuting $u,v$ generate $\widehat{\mathbb Z}^2$, the same
argument gives injectivity of $v-1$ on
\[
 \Lambda_P/(u-1)\Lambda_P
  =\widehat{\mathbb Z}[[\langle u\rangle\backslash P]].
\]
The $v$-orbits in this coset space have free procyclic action;
their finite lengths admit simultaneous refinements by any
prime power. The displayed quotient is the stated completed
permutation module, since the image of $u-1$ is compact and
therefore closed. Suppose now that a normal abelian subgroup
$D\subset P$ contains such $u,v$. For a continuous cocycle
$b:P\to\Lambda_P$, commutation gives
\[
 (v-1)b(u)=(u-1)b(v).
\]
The quotient injectivity first gives $b(u)=(u-1)t$; regularity
then gives $b(v)=(v-1)t$. Subtracting this coboundary, commutation
with $u$ forces $b|_D=0$, and normality forces every $b(g)$ to
be $D$-invariant and hence zero. We have proved
\[
 H^1_{\rm cts}(P,\Lambda_P)=0
 \tag{B.22.1}
\]
with compact continuous coefficients, without replacing them
by discrete modules. Apply this to $P=A$ and to $P=M$:
in the second case $c=[x,y]$ and $b=xcx^{-1}$ in $D=M'$
correspond to independent vectors $1,X$ in its regular module,
and generate a copy of $\widehat{\mathbb Z}^2$.

A class-preserving automorphism of $M$ acts identically on $A$
and sends $c$ to a monomial multiple $a c$. Remove an inner
automorphism lifting $a$. Its remaining difference from identity
is a cocycle from $A$ to its regular module, zero by (B.22.1).
Thus it is inner.

For level three put $L=F''/F'''$ and
$\Lambda=\widehat{\mathbb Z}[[M]]$. The completed Fox sequence is
\[
 0\longrightarrow L\longrightarrow\Lambda^2
 \xrightarrow{\partial}\Lambda
 \xrightarrow{\epsilon}\widehat{\mathbb Z}\longrightarrow0,
 \qquad\partial(a,b)=a(x-1)+b(y-1).
 \tag{B.22.2}
\]
This is the exact sequence, with these left-module conventions,
in \cite[Proposition 1.1]{PBC_YamaguchiCenters2026}; its source
is Ihara's Theorem A-2. Only this explicit sequence is used.
An $L$-valued cocycle bounds in $\Lambda^2$ by (B.22.1).
Its bounding vector $w$ has $(g-1)\partial w=0$ for all $g$;
regularity of $c-1$ forces $\partial w=0$. Hence
\[
                         H^1_{\rm cts}(M,L)=0.
 \tag{B.22.3}
\]

We need an explicit free cyclic submodule, including at every
coefficient prime. Let $r=[c,xcx^{-1}]\in L$, and put
$d=c-1$, $d^x=b-1$, $A_0=(c-1)x-(b-1)$.
The left Fox rule gives the literal identities
\[
 \delta_y(c)=x-c,\qquad
 \delta_y(r)=A_0(x-c),\qquad
                A_0d=d\,d^x(x-1).
 \tag{B.22.4}
\]
The elements $c,b,x,xc^{-1}$ have infinite primary order at
every prime. All indicated right multiplications are therefore
injective, even after coefficient reduction modulo $p$.
The last identity implies that $A_0$, and then $A_0(x-c)$,
has no left annihilator. Thus $r$ generates a free left
$\Lambda$-module, with the same assertion modulo every $p$.
No domain property of the full completed group ring is assumed.

Now remove the inner image on $M$ of a class-preserving
automorphism of $S_3$. Its action $T$ on $L$ is $\Lambda$-linear.
For any $v$ generating a free cyclic module, write $T(v)=gv$
by class preservation, with $g\in M$. Apply class preservation
also to $(1+2h)v$, for arbitrary $h\in M$. Freeness yields
\[
                  g+2hg=g_h+2g_hh.
\]
Distinct group elements are linearly independent over
$\widehat{\mathbb Z}$ in their completed group algebra.
For $h\ne1$ the two distinct terms on either side have distinct
coefficients $1,2$, so $g_h=g$ and $hg=gh$. Thus $g$ is central
in $M$, hence one. This argument includes coefficient prime two.

To reach any $l\in L$, choose the scalar
$\nu\in\widehat{\mathbb Z}$ whose $p$-coordinate is $p$.
The second Fox coordinate of $r+\nu l$ reduces modulo $p$ to
that of $r$. Successive division by $p$ proves its regularity
over $\mathbb Z_p[[M]]$, for every $p$; hence it also generates
a free cyclic $\Lambda$-module. Both $r$ and $r+\nu l$ are
fixed by $T$. Torsion-freeness of each coefficient module gives
$\nu(Tl-l)=0\Rightarrow Tl=l$.
The residual automorphism is identity on $M$ and $L$, so is a
continuous cocycle $M\to L$. Equation (B.22.3) makes it inner.
This proves the level-three rigidity and, by the first part with
$m=3$, the asserted fourth-to-third projection.
\end{proof}

\begin{statement}{B.23}{Proved}{Unrestricted arithmetic kernel classes and their exact peripheral failure}
For each odd prime $\ell$, the kernel of
$\Out(S_3)\to\Out(M)$ in the arithmetic centralizer contains
a closed copy of $\mathbb Z_\ell^\times$. Every nonidentity
member fails even to preserve the conjugacy class of
$\langle x\rangle$. More generally, a nonzero central Haar
idempotent supported on a normal pro-prime-to-$\ell$ subgroup
of $M$ cannot make any of $x-1,y-1,xy-1$ invertible.
These are the valid conclusions of LM6 A.59 and A.60;
they do not produce a peripheral element of the same-level kernel.
\end{statement}
\begin{proof}
Use $D=M'$, $L=F''/F'''$, and the Fox sequence (B.22.2).
The continuous Fox map identifies $S_3$ with
\[
 \{(v,g)\in\Lambda^2\rtimes M:\partial v=g-1\},
 \qquad\delta(fg)=\delta(f)+\overline f\delta(g).
 \tag{B.23.1}
\]
Indeed the sequence identifies its kernel over $M$ with $L$;
subtract the Fox vector of any lift of a prescribed $g$ to
obtain every vector having the indicated boundary.

The augmentation of $D=Rc$ modulo two gives an $M$-invariant
surjection $D^{(\ell')}\to C_2$ on its prime-to-$\ell$ factor.
Let $H$ be its kernel, and let $e_H$ be normalized Haar averaging
over $H$ in $\mathbb Z_\ell[[M]]$. Every finite denominator
is an $\ell$-adic unit. Choose $z=c^\eta$, where $\eta_\ell=0$
and $\eta_p=1$ for $p\ne\ell$, and put
\[
                         e=e_H(1-z)/2.
\]
This is a central idempotent of augmentation zero, regarded as
zero in the other coefficient factors. It is arithmetic-invariant:
the augmentation of the commutator module transforms by
$\chi^2$, which is one modulo two. Since $ez=-e$, the vector
$w=-e\delta([x,y]^\eta)/2$ satisfies $\partial w=e$.
For $t\in\mathbb Z_\ell^\times$ set $u_t=1+(t-1)e$ and
\[
 \Phi_t(v,g)=\bigl(u_t v+(1-t)(g-1)w,g\bigr).
 \tag{B.23.2}
\]
The cocycle identity for $(g-1)w$ and centrality of $u_t$ verify
the product and boundary equations directly. Moreover
$\Phi_t\Phi_s=\Phi_{ts}$, so these are continuous automorphisms.
They act identically on $M$ and by $t$ on $eL$.

For completeness $H^1_{\rm cts}(M,eL)=0$ has a short direct
proof in these exact compact coefficients. A cocycle restricted
to $H$ is a homomorphism from a pro-prime-to-$\ell$ group to a
pro-$\ell$ group, since $H$ acts trivially on $eL$. It is zero.
On $M/H$, the element $z$ is a central involution acting as
$-1$. Comparing the cocycle at $zg=gz$ gives
$2b(g)=(1-g)b(z)$, which is a coboundary because $\ell$ is odd.
The commutator of $\Phi_t$ with any arithmetic automorphism is
identity on $M,L$ and modulo $eL$. It therefore is such a
cocycle and is inner. Hence the outer classes centralize
arithmetic, over $\mathbb Q$ and every number field.

Here is an integral detector of their peripheral failure. Put
\[
                   n_x=e(1,0)-(x-1)w\in eL.
\]
Map $M$ to the dihedral group of order eight by the matrices
\[
 x\longmapsto\begin{pmatrix}1&0\\0&-1\end{pmatrix},\qquad
 y\longmapsto\begin{pmatrix}0&-1\\1&0\end{pmatrix}.
\]
The commutator is $-I$, $H$ maps to identity, and $e$ maps to
$I$. Apply the first row to the first Fox coordinate. It kills
$(x-1)eL$ because that row is fixed by $x$, but takes $n_x$
to $(1,0)$. Thus $(t-1)n_x\notin(x-1)eL$ for $t\ne1$.
Formula (B.23.2) changes $x$ by precisely this element.
If its image were conjugate to $x$, the conjugator in $M$
would belong to $C_M(x)=\langle x\rangle$; remove that power
to put the conjugator in $L$. Its conjugation difference lies
in $(x-1)L$, a contradiction after applying $e$.
Abelianization excludes a different exponent of $x$ as well.
The same detector shows that $t$ acts nontrivially on $eL$.
An inner map over center-free $M$ has conjugator in $L$ and
acts trivially on $L$, proving injectivity. Compactness of
$\mathbb Z_\ell^\times$ makes the image closed.

Finally let $N\triangleleft M$ be any pro-prime-to-$\ell$
subgroup. It commutes with the normal pro-$\ell$ factor
$D_\ell$ of $M'$, since their intersection is trivial.
The action on $D_\ell=\mathbb Z_\ell[[A]]c$ is the faithful
monomial action of $A$, so $N\subset M'$.
Consequently $N\cap\langle x\rangle=1$, and the continuous
map of profinite sets $N\to\langle x\rangle\backslash M$
is a closed embedding. It induces an injection of completed
free $\mathbb Z_\ell$-modules: every finite clopen partition
of a closed subset extends to one of the whole profinite space.
Therefore
\[
 \mathbb Z_\ell[[N]]\lhook\joinrel\longrightarrow
 \mathbb Z_\ell[[M]]/(x-1)\mathbb Z_\ell[[M]]
\]
is injective. A nonzero $e$ supported on $N$ cannot be in the
image of $x-1$, so $x-1$ cannot be onto $e\mathbb Z_\ell[[M]]$.
The identical argument applies to $y$ and $xy$, whose cyclic
subgroups also inject into $A$. This excludes the proposed
averaging method, not every conceivable peripheral construction.
\end{proof}

\begin{statement}{B.24}{Proved}{Independent finite certificates for the first two coefficient-five digits}
The explicit covers and point-count contradictions of LM6
A.54 and A.55 are correct. For the R28 unit family, a
third-derived arithmetic-centralizing lift over $\mathbb Q$
or $K_0=\mathbb Q(\mu_{15},E[5])$, where
$E:y^2=x^3+16$, forces $t_5\equiv1\pmod{25}$.
The complete exclusion at every prime and depth is B.20;
these finite certificates verify its earlier special cases.
\end{statement}
\begin{proof}
We give the geometric descent checks as well as the recomputed
integers. For $n=5$ or $25$, take the two covers
\[
 E\xrightarrow{[n]}E\xrightarrow{(y+4)/8}\mathbb P^1,
 \qquad X^n+Y^n+Z^n=0.
\]
For the second, put $r=X^n/Z^n$,
$w_2=(r-\omega)/(r-\omega^2)$ and
$t=w_2^3/(w_2^3-1)$, where $\omega^3=1\ne\omega$.
The respective cubic subcovers are
$w_1^3=t(t-1)$, with $w_1=x/4$, and $w_2^3=t/(t-1)$.
Their valuation vectors $(1,1)$ and $(1,-1)$ are independent
modulo three. The deck groups over
$K_n=\mathbb Q(\mu_{3n},E[n])$ are
\[
 G_n=(E[n]\rtimes C_3)\times((\mathbb Z/n)^2\rtimes C_3).
\]
Each factor has abelianization $C_3$, with fixed-point-free
order-three action on its $n$-torsion. A common nontrivial
quotient would be solvable and have a nontrivial abelian quotient,
forcing a common cubic subfield. The independence excludes this;
thus the fibre product is connected. The same action shows
$Z(G_n)=1$.

The R28 unit acts as $(t_5,1)$ on the two derived modules.
Extend these scalars by identity on the complement. Since
$q(x)-1$ is invertible on both modules, a unique inner correction
fixes $q(x)$. Equality on the commutator then forces equality
on $q(y)$, again by that invertibility. Hence this is an actual
induced automorphism of the finite quotient, not just an
endomorphism of its derived module. The inner correction does
not change its action on the latter.

Here are all parameters and the independently reproduced counts:
\[
\begin{array}{c|cc}
 &n=5&n=25\\\hline
 q&211&5701\\
 \omega&14&75\\
 \zeta_n&55&303\\
 P_0&(103,117)&(25,2322)\\
 v_0&(3,1)&(0,23)\\
 |E(\mathbb F_q)|&225&5625\\
 |E[n](\mathbb F_q)|&25&625\\
 \text{exponent of }E(\mathbb F_q)&15&75
\end{array}
\]
In particular these primes split completely in $K_n$.
Set $J_a=\langle(aP_0,v_0)\rangle$. For each rational
$t\ne0,1$ having both cubic classes trivial, choose their cube
roots $w_1,w_2$ and put
\[
 Q=(4w_1,8t-4),\quad
 r=\frac{w_2\omega^2-\omega}{w_2-1},\quad
 v=\left(\log_{\zeta_n}r^{(q-1)/n},
       \log_{\zeta_n}(-1-r)^{(q-1)/n}\right).
\]
These logarithms are the unique finite exponents modulo $n$.
The elliptic Kummer vector is conjugate under $C_3$ to $3Q$
when $n=5$, and to $-3Q$ when $n=25$.
Indeed the Frobenius trace and norm in $\mathbb Z[\omega]$
give respectively $1+15\omega^i$ and $1-75\omega^i$;
applying $(\pi-1)$ to an $n$-division point proves these
formulas. Thus, writing this vector as $u$ and
$T(a,b)=(b-a,-a)$, the exact fixed-coset count is
\[
       n^3\sum_{i,j=0}^{2}
            \mathbf1_{J_a}(\omega^i u,T^jv).
 \tag{B.24.1}
\]
For each complement coordinate there are $n^3$ translation
cosets; the subgroup-membership condition is precisely the
condition that they are fixed. This proves the formula directly.
Base values with nontrivial cubic image contribute zero.

The program \texttt{pbc\_lm6\_counts.py} was written
independently, with its own elliptic group law and subgroup
membership sets; it imports no supplied certificate code.
It checks all base values and all nine root pairs at each
eligible value. There are $24$ eligible values and $216$ root
pairs for $n=5$, giving
\[
 \#(Y_5/J_1)(\mathbb F_{211})=250,
 \qquad\#(Y_5/J_4)(\mathbb F_{211})=0.
\]
The only contributions to the first count are $125$ each at
$t=12,200$. For $n=25$ it checks $624$ eligible values and
$5616$ root pairs, giving
\[
\begin{array}{c|rrrrr}
 a&1&6&11&16&21\\\hline
 \#(Y_{25}/J_a)(\mathbb F_{5701})&15625&0&0&0&15625.
\end{array}
\]
The sole contributing fibres are $t=8$ for $a=1$ and
$t=4104$ for $a=21$. The complete executed output is in the
adjacent \texttt{.out} and \texttt{.json} files.

A hypothetical arithmetic-centralizing lift induces the
arithmetic quotient action exactly: the deck group is constant
and center-free, so its possible central arithmetic twist is
zero. The component subgroups contain $F''$ and their derived
groups contain $F'''$, retaining their full first cohomology at
level three. These explicit marked covers have good tame
reduction at the displayed primes. The duality and trace
comparison of B.20 forces equal point counts for subgroups
carried to each other. If $t_5\not\equiv1\pmod5$, a power of
its residue is $-1$, taking $J_1$ to $J_4$ and contradicting
the first comparison. If it is one modulo five but not modulo
twenty-five, a power is six modulo twenty-five, taking $J_1$
to $J_6$ and contradicting the second. This proves the claimed
second-digit restriction and completes the independent audit.
\end{proof}

\subsection{Relative levels and the precise remaining involution problem}

\begin{statement}{B.25}{Proved}{The inverse-system and odd-prime conjugacy reductions in Lead 26(b)}
For open characteristic subgroups $U\subset F$, put
$N_U=R_{\mathrm{nil}}(U)$ and $Q_U=F/N_U$.
The natural homomorphisms are onto and give
\[
 \mathcal P\simeq
 \varprojlim_U\Out^\sharp(Q_U).
 \tag{B.25.1}
\]
Here $\Out^\sharp$ preserves the three specified peripheral
classes up to powers. Involutions in $\mathcal P$ are precisely
compatible families of elements of order dividing two in this
system, nontrivial at some level. Two such involutions are
conjugate in $\mathcal P$ exactly when their images are
conjugate at every level.

Writing $D=F/U$ and $O^p(U)=\ker(U\to U^{(p)})$, one has
\[
 Q_U\simeq\mathop{\prod}_{p,D}F/O^p(U).
 \tag{B.25.2}
\]
In a component $E_p=F/O^p(U)$, consider automorphisms
preserving its marked quotient $D$ and its kernel $U^{(p)}$.
The kernel on outer classes of reduction to $F/(U'U^p)$ is
pro-$p$. At odd $p$, any two involutions with the same image
in this finite outer quotient are conjugate by that kernel.
This applies in particular to the component automorphisms
induced from $\Out^\sharp(Q_U)$.
\end{statement}
\begin{proof}
First
\[
                    N_U=[N_U,U]=[N_U,F],
 \tag{B.25.3}
\]
where commutators mean closed subgroups. Indeed the image of
$N_U$ in $U/[N_U,U]$ is central and its quotient is the
pronilpotent group $U/N_U$. Every finite quotient of this
central extension is nilpotent. Thus $U/[N_U,U]$ itself is
pronilpotent, and the defining property of $N_U$ gives the
first equality. Normality in $F$ gives the second.
B.11 now lifts every multiplier-one peripheral automorphism
of $Q_U$ with its three frames, since its central lifting
quotient is $Q_U$ itself. For an arbitrary multiplier, compose
with the inverse of a Galois automorphism having that
cyclotomic multiplier, lift, and compose back. Such an
automorphism preserves $U$ and $N_U$ by characteristicity.
The three exponents in the definition of $\Out^\sharp$ must
be a common unit: $N_U\subset F'$ retains the independent
abelian vectors of $x,y$, while $z$ has vector $-x-y$.
This proves all asserted surjectivities.

If $U\subset V$, then $N_U\subset N_V$, by restricting the
nilpotent quotients of $V$ to $U$. Every peripheral automorphism
of $Q_U$ lifts to $F$, so it preserves $N_V/N_U$ and defines
the transition homomorphism to $\Out^\sharp(Q_V)$. This also
explains why those transitions are well-defined even if one
does not identify the larger kernel intrinsically in an
arbitrary quotient automorphism group.
The intersection of the $N_U$ is trivial, since $N_U\subset U$
and characteristic open subgroups of finitely generated $F$
are cofinal. Hence $F=\varprojlim Q_U$.

The group $\mathcal P$ is compact: continuous endomorphisms
of $F$ form $F^2$, surjectivity is the closed condition of
generating every finite quotient, and surjective endomorphisms
of finitely generated profinite groups are automorphisms.
The peripheral conditions are closed, with compact sets of
multipliers and conjugators; quotienting by inner automorphisms
preserves compactness. If an element is inner on every $Q_U$,
the corresponding closed sets of conjugators in $F$ have the
finite intersection property, so it is inner on $F$. This gives
injectivity in (B.25.1). Surjectivity follows by the same
compactness argument: any finite compatible collection is the
image of an element at a common finer level and thus of an
element of $\mathcal P$. The statements about order follow
from this group isomorphism. For conjugacy use the nonempty
closed sets of conjugators in $\mathcal P$ at each level;
surjectivity lifts each level conjugator, and a common finer
level gives the finite intersection property.

The maximal pronilpotent quotient of $U$ is
$\prod_p U^{(p)}$. This follows directly from finite nilpotent
groups being the products of their Sylow groups: maps onto
finitely many different prime quotients are jointly onto,
since their orders are pairwise coprime. Compactness gives
the full product statement. Comparing kernels over $D$
proves (B.25.2).

We supply the kernel argument for its marked component groups.
Put $V=U^{(p)}\triangleleft E_p$; it is free pro-$p$ of rank
$1+[F:U]$. One has
\[
                         C_{E_p}(V)=1.
 \tag{B.25.4}
\]
To see this, the usual finite covering graph of the two-petal
rose, with deck group $D$, gives over $\mathbb Q_p$
\[
 0\longrightarrow V^{\rm ab}\otimes\mathbb Q_p
  \longrightarrow\mathbb Q_p[D]^2
  \longrightarrow\mathbb Q_p[D]
  \longrightarrow\mathbb Q_p\longrightarrow0.
\]
Semisimplicity in characteristic zero gives
$V^{\rm ab}\otimes\mathbb Q_p\simeq
\mathbb Q_p\oplus\mathbb Q_p[D]$ as $D$-modules. This is
faithful, even when $p$ divides $|D|$. Thus an element
centralizing $V$ has trivial image in $D$ and lies in $Z(V)$,
which is trivial because the free pro-$p$ rank is at least two.

Let $K_p$ be the kernel of the marked automorphism action on
$E_p/\Phi(V)=F/(U'U^p)$, where
$\Phi(V)=\overline{V^p[V,V]}$. Restriction embeds $K_p$ into
$\ker(\Aut(V)\to\Aut(V/\Phi(V)))$: if an automorphism fixes
$V$ pointwise, its discrepancy at any $g\in E_p$ centralizes
$V$, and (B.25.4) makes it trivial. The latter automorphism
kernel is pro-$p$. For a direct verification, use the lower
$p$-central filtration of the finitely generated free pro-$p$
group. An automorphism trivial on the first quotient is trivial
on each associated commutator-and-power graded piece, and its
successive congruence kernels embed in finite groups of linear
maps between finite $\mathbb F_p$-vector spaces. Thus all its
finite congruence quotients are $p$-groups, and the separated
filtration gives the asserted pro-$p$ inverse limit.
Consequently $K_p$ is pro-$p$.

If an outer class reduces to an inner class in the finite
quotient, remove an inner representative lifted from $E_p$.
It is then represented in $K_p$. The outer kernel is therefore
a quotient of this compact pro-$p$ group and is pro-$p$.
For odd $p$, two involutions with the same residue are two
complements to this kernel in its extension by $C_2$.
They are conjugate by the kernel. This coprime-complement
fact can be seen in finite $p$-group quotients, inducting along
a central series and using that $H^1(C_2,A)=0$ for each finite
abelian $p$-group $A$ because multiplication by two is invertible;
compactness then chooses one compatible conjugator.

All component automorphisms coming from $Q_U$ preserve the
marked $U$ and $D$, because they lift to $F$. If arbitrary
component automorphisms are considered instead, this marking
condition is part of the reduction map, not an implicit claim
that every characteristic subgroup of $F$ remains intrinsic
in every quotient. A cofinal system of fully invariant open
$U$ removes this caveat: intersect kernels of all maps to
finite groups of bounded order; every component endomorphism
lifts to an endomorphism of free $F$ and preserves such a $U$.
\end{proof}

\begin{statement}{B.26}{Refuted}{Parity is not a character of the first unframed relative level}
For $H=F'F^3$, the character $\Delta$ does not factor through
the image of $\widehat{\mathrm{GT}}_0$ in
$\Out^\sharp(F/R_{\mathrm{nil}}(H))$.
In particular, assigning a $\Delta$ value to every unframed
conjugacy class at that relative level is not well-defined.
This corrects the parity-reading suggestion in Lead 26(b).
It does not decide which of the two parity values can occur
among full involutions lifting one specified level involution.
\end{statement}
\begin{proof}
The exact E.58 element in R24 has $f\in R_{\mathrm{nil}}(H)$
and $\Delta(1,f)=1$. Its normalized automorphism is
$x\mapsto x$, $y\mapsto f^{-1}yf$, so it induces identity
on $F/R_{\mathrm{nil}}(H)$. The identity full element has the
same relative image and parity zero. This proves nonfactorization.
Indeed multiplication by the E.58 element toggles parity in
every fibre of the group-image map. Such multiplication need
not preserve the involution equation, so it gives no inference
about the parity of involutive lifts.

The distinction concerns frames, not a failure of the product
isomorphism. E.58 retains a nontrivial two-primary $r$-frame
for an automorphism which is identity on the entire relative
source. Uniqueness of the pure torsion frames holds in the full
modular group; it fails in the relative quotient, where those
frames may lie in enlarged torsion centralizers. Consequently
forgetting the frame loses information that $\Delta$ uses.
The requested $U_{12}$ class analysis, and the compatible
involutive lift analysis at either level, remain open.
\end{proof}

\begin{statement}{B.27}{Proved}{The cubic metabelian involution has a full gentle lift}
The metabelian involution $\varphi_0$ is the image of an
element of the full profinite group
$\widehat{\mathrm{GT}}_0$ of multiplier one. Every such
full lift has infinite outer order and lies outside
$\widehat{\mathrm{GT}}$.
More generally, the entire cubic sign family
$\prod_{p\geq5}C_2$ in R28 is contained in the metabelian
image of the full gentle group. This strengthens the
$\mathcal P_1$ lifting conclusion of B.3 and supplies a
second proof of strict full profinite containment in P.4.
\end{statement}
\begin{proof}
We keep the actual modular source throughout. Write
\[
 \Gamma=C_2\amalg C_3=\langle s,r\mid s^2=r^3=1\rangle,
 \quad\tau=sr,\quad F=\ker(\Gamma\to S_3),\quad
 N=F'',\quad\overline\Gamma=\Gamma/N.
\]
Use $x=\tau^{-2}$, $y=sxs$, as in P.4. The metabelian
unit $\beta_0$ is unchanged by that specified interchange
and inversion of the original marking.
The metabelian
two-cycle and three-cycle identities of B.9 give an actual
automorphism $\alpha_0$ of $\overline\Gamma$ by
\[
             s\mapsto sf_0,\qquad r\mapsto f_0^{-1}r,
                         \tau\mapsto\tau.
 \tag{B.27.1}
\]
Its restriction to $M=F/N$ is the normalized involution
$\varphi_0$. The centralizer of $M$ in $\overline\Gamma$
is trivial: its intersection with $M$ is $Z(M)=1$, and
the faithful $S_3$ action on $M^{\rm ab}$ detects every
remaining coset. Thus $\alpha_0^2=1$ on the whole extension.
The element $f_0$ belongs to the pro-five factor
$V=M'_5$. Therefore $\alpha_0$ is identity on
$\Gamma_0=\overline\Gamma/V$.
Moreover its images of $s,r$ are conjugate to $s,r$ by
elements $b_s,b_r\in V$: in each preimage of the same
cyclic subgroup of order two or three in $\Gamma_0$, they
are complements to the normal pro-five group. Coprime
complement conjugacy gives the stated conjugators.
The third frame is $b_\tau=1$.

Consider the actual centralization of this specified kernel,
\[
 E=\Gamma/[N,\Gamma],\qquad
 K=N/[N,\Gamma],\qquad E_q=E/\prod_{p\ne q}K_p.
 \tag{B.27.2}
\]
All subgroups are closed, and $K_p$ denotes the pro-$p$
factor of the abelian group $K$. We will construct a
peripheral automorphism of $E$ above (B.27.1).
The homology of a profinite free product is the sum of
the homologies of its finite factors in degrees at least
two. For a finite cyclic group with trivial
$\mathbb Z_q$ coefficients, its periodic resolution gives
$H_2(C_n,\mathbb Z_q)=0$: the relevant kernel is that
of multiplication by $n$ on torsion-free $\mathbb Z_q$.
Thus $H_2(\Gamma,\mathbb Z_q)=0$ for every $q$.
Since $N\subset\Gamma'$ and abelianization is unchanged,
the continuous five-term sequence gives the natural
identification
\[
                  K_q\simeq H_2(\overline\Gamma,\mathbb Z_q).
 \tag{B.27.3}
\]
This is a Hopf calculation for the source and kernel in
(B.27.2); no assertion about a universal central extension
of arbitrary profinite groups is being used.

At $q=5$, quotienting first by $[N,F]$ also gives
\[
       K_5=(H_2(M,\mathbb Z_5))_{S_3}=(C_5)_{S_3}.
\]
The quotient is the continuous coinvariant module. Since
six is invertible, averaging identifies it with a direct
summand of the torsion-free module $C_5$ in B.12. The action
of $\alpha_0$ on $C_5$ is multiplication by
$\beta_0\beta_0^*=1$, by the completed exterior calculation
there. Its action on $K_5$ is consequently identity.

Every automorphism of $\overline\Gamma$ has a unique lift
to an automorphism of $E_5$. To establish existence,
lift its images of $s,r$ to their unique elements of orders
two and three in the corresponding central fibres.
These unique torsion lifts exist because multiplication
by two and three is invertible on $K_5$.
The free-product presentation defines a homomorphism
$\Gamma\to E_5$. It sends $N$ into the central group
$K_5$, kills $[N,\Gamma]$ and its non-five coefficient
factors, and hence factors through $E_5$.
Its action on $K_5$ is the functorial action in (B.27.3):
naturality of transgression compares this map with the
original source map $\Gamma\to E_5$. It is an isomorphism
on $K_5$ because that original quotient identifies the target
transgression $H_2(\overline\Gamma,\mathbb Z_5)\to K_5$
with (B.27.3); the new map therefore acts on $K_5$ by
$H_2(\alpha_0,\mathbb Z_5)$.
It is an isomorphism
on both kernel and quotient, and therefore is an
automorphism. For uniqueness, an automorphism over identity
on $\overline\Gamma$ acts identically on $K_5$ by the same
Hopf argument. Its difference from identity is a
homomorphism $\overline\Gamma\to K_5$, which is zero
because $\overline\Gamma^{\rm ab}=C_6$.

Let $\widetilde\alpha_5$ be this unique lift of
$\alpha_0$. Uniqueness gives
$\widetilde\alpha_5^2=1$. It acts identically on $K_5$,
and $\widetilde\alpha_5(\tau)=\tau a$ for some
$a\in K_5$. Squaring gives $2a=0$, hence $a=0$.
For either prescribed frame $b_s,b_r$, choose any lift
to $E_5$. Conjugating $s$ or $r$ by that lift gives
the unique torsion lift above its specified image, so
these are peripheral frames for $\widetilde\alpha_5$.
Thus all three peripheral conditions, with the retained
frame coordinates, hold at five.

For $q\ne5$ we give the exact comparison of central
extensions. Let $N_0$ be the inverse image of $V$ in
$\Gamma$ and define
$E_{0,q}=\Gamma/[N_0,\Gamma]$ with its non-$q$ central
factors killed. Again $N_0\subset\Gamma'$, so its
central kernel is $H_2(\Gamma_0,\mathbb Z_q)$ by
the same five-term calculation. Averaging over the
pro-five normal subgroup $V$ gives
$H_i(V,\mathbb Z_q)=0$ for $i>0$. The continuous
Hochschild--Serre sequence therefore gives an isomorphism
\[
 H_2(\overline\Gamma,\mathbb Z_q)
                 \xrightarrow{\sim}H_2(\Gamma_0,\mathbb Z_q).
\]
The natural map $E_q\to E_{0,q}$ consequently induces
an isomorphism on their central kernels. Comparing the
two central exact sequences over $\overline\Gamma$
proves the Cartesian square, equivalently
\[
                  E_q\simeq
             \overline\Gamma\mathbin\times_{\Gamma_0}E_{0,q}.
 \tag{B.27.4}
\]
The kernel and quotient maps are isomorphisms, so this
conclusion follows by a direct short-five-lemma argument
for compact groups, not by a choice of splitting.

On (B.27.4) use the automorphism
$(\alpha_0,\operatorname{id})$. It fixes $\tau$
literally. The pairs $(b_s,1),(b_r,1)$ are elements
of this fibre product because $b_s,b_r\in V$.
They are the required conjugators of $s,r$.
This proves all peripheral conditions at every other
coefficient prime, including two and three.

The central decomposition identifies
$E=\prod_{q,\overline\Gamma}E_q$. The constructed maps
and frames have the same base coordinates, so assemble
to a peripheral automorphism of $E$ above $\alpha_0$
with third frame one. The proof of B.11 applies without
change to the modular source: the finite scalar-cover
and counting argument in P.3 already enforces the
relations $s^2=r^3=1$ and uses only normality of the
finite image of $N$. It therefore lifts this automorphism
to the full $\Gamma$ with those frame coordinates.
An inner correction in the kernel makes $\tau$ fixed
literally. Set $f=s\Phi(s)$. The retained image is
$f_0\in M'$, so $f\in F'$, and the two torsion
relations give the full gentle two-cycle and three-cycle
identities exactly as in P.4. The restriction to $F$
is an automorphism. Thus $(1,f)\in
\widehat{\mathrm{GT}}_0$ has metabelian image
$\varphi_0$.

All arguments work with any one coefficient prime
$p\geq5$ in place of five for the cubic sign involution.
Finite products of the resulting full gentle elements
realize every finite set of signs in R28. Their
metabelian images converge, coordinatewise, to any
specified infinite sign tuple. The image of the compact
full gentle group in $\Out(M)$ is closed, so it
contains the entire product $\prod_{p\geq5}C_2$.
The finite-order exclusion is still R25/B.13: no
nontrivial full multiplier-one lift has finite order.
Finally B.2 and B.9 show that the beta two-cocycle
identity forced by the pentagon fails for $\varphi_0$.
Every full gentle lift just constructed therefore lies
outside the full pentagon group, proving the strictness
asserted here independently of the degree-seven route.
\end{proof}

\begin{statement}{B.28}{Proved}{The negative residue survives the finite detector and is Galois-realized}
Put $\tau_0=\varphi_0c$ in $\Out(M)$, where $c$ is the
normalized complex conjugation. This is an involution of
multiplier $-1$. On the quotient $G$ of order $614400$ in
B.19, exactly four of the twenty possible involutions above
its prescribed $Q_{75}$ action satisfy both negative
peripheral frame equations and both frame norm equations.
Moreover the specified action of $\tau_0$ on $Q_{75}$ is
already the action of a genuine Galois involution.
These are finite survival statements, not a construction
of an involutive full lift of $\tau_0$.
\end{statement}
\begin{proof}
The unit $\beta_0$ is invariant under $*$ and has square
one, so $\varphi_0$ commutes with $c$ and $\tau_0^2=1$.
Keep all coordinates, the field $\mathbb F_{16}$, and the
pair $x,y$ of B.19. The action on the quotient is
\[
 (a,b,i)\longmapsto(-a,b-a,-i).
 \tag{B.28.1}
\]
The three characters defining $W=\mathbb F_{16}^3$ are
preserved as a set: the involution $C(a,b)=(a,a-b)$
fixes the first and interchanges the other two, and the
additional negative sign is fourth-power Frobenius.
Thus the multiplicity-one descent argument in B.19
applies to this action as well. Its intertwiners are
\[
        L_a(u)_j=a u_{-j}^{,4},\qquad a\in\mu_5.
\]
The central correction is $Q_0(u)$, exactly as in
B.19.5. All lifts are $\operatorname{Inn}(k)\psi_a$.
The conditions on $x$ and on the square are
$k=(b,b,b)$ and $ab^4=b$, leaving twenty candidates.

The required quotient coordinates of the frames are
$q_f=(3,4,0)$ and $q_h=(4,3,1)$. The independent
enumeration in \texttt{pbc\_extraspecial\_negative.py}
tests
\[
 \alpha(y)=f^{-1}y^{-1}f,\quad
 \alpha(xy)=h^{-1}(xy)^{-1}h,\quad
 f\alpha(f)=h\alpha(h)=1
 \tag{B.28.2}
\]
for every candidate and every possible frame in those
fibres. All twenty have eight $f$ choices. Exactly four
also have eight $h$ choices: $a=\zeta^2$ and the packed
coordinates $k=0,1365,2184,3549$. For $k=0$, frames with
central bit zero are
\[
       f_W=(0,\zeta^3,1),\qquad h_W=(1,\zeta,0).
\]
The program verifies the field law, the group cocycle,
the intertwining identities, all twenty automorphisms,
and every equation in (B.28.2). It writes the complete
table to \texttt{pbc\_extraspecial\_negative.json}.
These four surviving automorphisms are mutually conjugate
by diagonal inner elements of the extraspecial group:
on the diagonal copy of $\mathbb F_{16}$, the image of
$1+L_{\zeta^2}$ is its fixed $\mathbb F_4$ line, which
is exactly the set of four allowed $k$.

There is also an arithmetic reason why the order-$75$
residue alone cannot yield an obstruction. Use the elliptic
curve $E:y^2=x^3+16$ and $t=(y+4)/8$ of B.24.
Its degree-$75$ cover is the quotient after $[5]$, with
deck group $E[5]\rtimes C_3$. At $13$ direct enumeration
gives $\#E(\mathbb F_{13})=9$, so the Frobenius trace
is $5$ and its determinant on $E[5]$ is $13\bmod5$.
The trace formula and determinant statement are
\cite[V, Proposition 8.3 and the Weil-pairing discussion
after Conjecture 8.5]{PBC_MilneEC2021}.
Since $13\equiv1\pmod3$, this Frobenius $d$ commutes
with the cubic rotation and belongs to the nonsplit
Cartan $\mathbb F_{25}^{\times}$. Its trace is zero
modulo five, hence $\bar d=-d$. Complex conjugation
acts on the Cartan by $d\mapsto\bar d$.

Choose a Galois element $\sigma$ above this Frobenius and
use the normalized representative fixing $x$ on $Q_{75}$;
this is possible because its multiplier is one modulo
three. Then $\sigma c\sigma^{-1}$ acts on $E[5]$ by
$(d/\bar d)C=-C$ and sends the cubic rotation to its
inverse. It therefore gives precisely (B.28.1).
The exact point count and both possible Frobenius
matrices are checked in
\texttt{pbc\_negative\_shadow.py}. A genuine full Galois
involution consequently has this order-$75$ action.
Any obstruction to a full involutive lift of the entire
metabelian $\tau_0$ must retain additional metabelian
information.
\end{proof}

\begin{statement}{B.29}{Proved}{A negative metabelian involution not conjugate to complex conjugation}
Let $\mathcal I\subset\Out(M)$ be the metabelian image
of the full pure peripheral group $\mathcal P$.
The involution $\tau_0=\varphi_0c$ belongs to $\mathcal I$
and is not conjugate to $c$ by any element of $\mathcal I$.
In fact it lies in the metabelian image of
$\widehat{\mathrm{GT}}_0$ by B.27. Thus any full
involutive lift of $\tau_0$, if one exists, is not
conjugate to complex conjugation even in $\mathcal P$.
Existence of such an involutive lift is not asserted.
\end{statement}
\begin{proof}
We use two specified metabelian quotients, not a chosen
irreducible quotient whose invariance would be uncertain.
First impose the triangle relations
\[
 (x^3,y^3,(xy)^3)=1,
 \qquad\text{or respectively}\qquad
 (x^4,y^4,(xy)^2)=1.
 \tag{B.29.1}
\]
In each case, take the kernel of the cyclic character
$x,y\mapsto1$ in $C_3$ or $C_4$, then quotient its
translation lattice by five. The resulting groups are
\[
 Q_3=\mathbb F_5^2\rtimes C_3=Q_{75},\qquad
 Q_4=\mathbb F_5^2\rtimes C_4,
 \tag{B.29.2}
\]
of orders $75$ and $100$. To check the translation
description directly, the cyclic covers
$w^3=t(t-1)$ and $w^4=t(t-1)$ have genus one by
Riemann--Hurwitz, with ramification signatures
$(3,3,3)$ and $(4,4,2)$. Their compact fundamental
groups are $\widehat{\mathbb Z}^2$. The cyclic deck
groups act by rotations of orders three and four.

Every pure peripheral automorphism preserves the
normal closures in (B.29.1), since it raises each
peripheral generator to a unit power up to conjugacy.
It also preserves the indicated cyclic kernel, by
scalar action on abelianization, and its subgroup of
five multiples. Thus both finite quotients are
preserved. Write $H_3,H_4$ for its two linear actions
on the translation modules. They normalize their
respective rotation groups, with exponent $\lambda$.

For any element of $\mathcal I$ the determinants obey
\[
                    \det H_3=\det H_4=\lambda_5.
 \tag{B.29.3}
\]
Indeed choose a Galois element of the same global
multiplier, using the surjectivity of the cyclotomic
character. Dividing by it leaves an element of
$\mathcal P_1$. By B.17 its metabelian unit satisfies
$\beta\beta^*=1$. On the two compact genus-one
modules the character pairs are respectively
$(X,Y)=(\omega,\omega),(\omega^{-1},\omega^{-1})$
and $(i,i),(-i,-i)$. The determinant of this
multiplier-one element is therefore one on both.
For the Galois element, (B.29.3) is the equivariance
of the Weil pairing on the two genus-one covers
defined over $\mathbb Q$; replacing a genus-one curve
by its Jacobian leaves this homology action unchanged.
See \cite[V, \S8, p.~221]{PBC_MilneEC2021}.
Inner translations act trivially on the translation
module and inner rotations have determinant one,
so these determinant statements do not depend on
the auxiliary arithmetic base points.

Now suppose that $hc h^{-1}=\tau_0$ in $\Out(M)$,
with $h\in\mathcal I$. By the group section in B.7
this is a literal equality of normalized automorphisms.
In particular no extra inner rotation can be inserted
independently in either quotient. The action of
$\varphi_0$ on the cubic translation module is $-I$;
on the quartic module it is $I$, since its coefficient
support is the specified cubic character block.
Consequently
\[
             H_3C_3H_3^{-1}=-C_3,\qquad
             H_4C_4H_4^{-1}=C_4,
 \tag{B.29.4}
\]
where $C_j$ is the action of normalized complex
conjugation, with $C_j^2=1$ and determinant $-1$.

For the cubic rotation its normalizer is
$\mathbb F_{25}^{\times}\rtimes\langle C_3\rangle$.
Write $H_3=dC_3^\epsilon$. The first equation in
(B.29.4) is $d/\bar d=-1$, so $d$ is a nonzero
trace-zero element. Its square is a nonsquare in
$\mathbb F_5$, and
$N(d)=-d^2$ is also a nonsquare because $-1$ is
a square in $\mathbb F_5$. Since
$\det C_3=-1$ is square, $\det H_3$ is nonsquare.

The quartic rotation splits over $\mathbb F_5$ with
eigenvalues $i=2$ and $-i=3$. In this eigenbasis
$C_4$ exchanges the two eigenlines. A normalizing
matrix is diagonal or antidiagonal. Commuting with
$C_4$ forces it to be $aI$ or $aC_4$, respectively.
Its determinant is $a^2$ or $-a^2$, both squares in
$\mathbb F_5$. This contradicts (B.29.3).

For an independent finite check,
\texttt{pbc\_negative\_shadow.py} enumerates all
$480$ matrices of $\mathrm{GL}_2(\mathbb F_5)$.
The cubic and quartic normalizers have orders $48$
and $32$. Exactly eight matrices satisfy each
respective equation of (B.29.4); their determinant
sets are $\{2,3\}$ and $\{1,4\}$, disjoint as
the proof requires. B.27 places $\tau_0$ in the
full gentle metabelian image, completing the claim.
\end{proof}

\begin{statement}{B.30}{Proved}{The coefficient-two cubic sign also has a full gentle lift}
Let $e$ be the idempotent in $R_2=\mathbb Z_2[[A]]$
supported on the two prime-to-two characters
$(X,Y)=(\omega,\omega)$ and $(\omega^2,\omega^2)$,
where $\omega$ is a primitive cube root of unity.
The metabelian unit which is $1-2e$ at coefficient
prime two and one at every other coefficient prime
is the image of an actual multiplier-one element of
$\widehat{\mathrm{GT}}_0$. Consequently the entire
cubic sign family $\prod_{p\ne3}C_2$ has full gentle
lifts. Each nonidentity member has only infinite-order
full lifts and has no pentagon lift.
\end{statement}
\begin{proof}
We refine B.27 at the one place where averaging over
$S_3$ was unavailable. Keep its notation
$N=F''$, $\overline\Gamma=\Gamma/N$,
$E=\Gamma/[N,\Gamma]$ and $K=N/[N,\Gamma]$.
Let $V=eM'_2$, let $N_0$ be its inverse image in
$\Gamma$, and put
\[
 \Gamma_0=\overline\Gamma/V,\qquad
 E_0=\Gamma/[N_0,\Gamma],\qquad
 K_0=N_0/[N_0,\Gamma].
 \tag{B.30.1}
\]
All these are actual source quotients. Both $N$ and
$N_0$ lie in $\Gamma'$, so the same Hopf calculation
as B.27 identifies their central kernels with the
respective second homology groups. A subscript two
on $E$ or $E_0$ means that the non-two central factors
have been killed.

We first compute precisely the kernel on the central
modules. Write $C_2=H_2(M,\mathbb Z_2)$ as in
B.12. The invariant idempotent $e$ acts on $C_2$:
indeed the additive submodule $J_2=R_2^\iota$ in
(B.12.1) is preserved by multiplication by $e$, since
$e^\iota=e$. The geometric $S_3$ action on this
description is the ordinary substitution action on
$R_2$. For the transposition the commutator generator
changes by a minus sign; for the three-cycle it changes
by a monomial. Their two exterior factors cancel the
sign or the monomial and its inverse. Thus this action
commutes with $e$.

For clarity the exact homology calculation after
killing $V$ is
\[
 H_2(M/V,\mathbb Z_2)\simeq(1-e)C_2,
 \qquad
 \ker(K_2\longrightarrow K_{0,2})
             \simeq(eC_2)_{S_3}=:L.
 \tag{B.30.2}
\]
To verify it, $M'_2=V\oplus W$ with
$W=(1-e)R_2$. The action of $X-1$ is invertible
on $V$, so $H_i(A,V)=0$ in every degree, including
zero. The regular-module calculation of B.12 then
gives $H_i(A,W)=0$ for $i>0$ and
$H_0(A,W)=\mathbb Z_2$. The boundary transgression
remains the isomorphism generated by $[x,y]$.
The same degree-two spectral sequence therefore
identifies $H_2(M/V,\mathbb Z_2)$ with
$(\widehat\bigwedge^2 W)_A$.
The cross summand $(V\widehat\otimes W)_A$ is zero:
after an unramified coefficient extension, diagonal
coinvariants pair only inverse prime-to-two
characters, and neither inverse of a character in
$V$ occurs in $W$. This can also be checked at every
finite character quotient and passed to the compact
inverse limit. The remaining exterior summands give
the first formula in (B.30.2).
Finally $N/[N,\Gamma]=(N/[N,F])_{S_3}$, and the
identical formula holds for $N_0$. Taking coinvariants
of the $S_3$-stable direct decomposition by $e$ proves
the second formula, including its surjectivity.

The crucial point is that $L$ is torsion-free, even
though we do not assert this for all of $K_2$.
Let $\mathcal O=\mathbb Z_2[\omega]$ and let
$A_2$ be the pro-two factor of $A$. Finite unramified
base change splits the two cubic character blocks:
\[
 eR_2\widehat\otimes_{\mathbb Z_2}\mathcal O
       \simeq\mathcal O[[A_2]]\oplus\mathcal O[[A_2]].
\]
Inversion exchanges the two summands and inverts
$A_2$. Quotienting by its invariant graph gives
\[
 eC_2\widehat\otimes_{\mathbb Z_2}\mathcal O
       \simeq\mathcal O[[A_2]].
 \tag{B.30.3}
\]
Each cubic character has triple $(\omega,\omega,\omega)$
or its inverse, so every geometric permutation fixes
that character. Its residual action in (B.30.3) is
exactly the permutation of $A_2$ induced by $S_3$.
Completed permutation-module coinvariants therefore
give
\[
 L\widehat\otimes_{\mathbb Z_2}\mathcal O
                \simeq\mathcal O[[A_2/S_3]].
 \tag{B.30.4}
\]
Here the compact orbit space is meant literally.
The equality follows from the universal property of
the completed free module on a profinite space:
quotienting by all orbit differences is the completed
free module on that orbit space. It does not use
division by six. This module is torsion-free; finite
faithfully flat base change proves the same for $L$.

Let $\alpha_0$ be the involution of $\overline\Gamma$
given by the two metabelian cycle identities in B.9.
It is identity on $\Gamma_0$, acts identically on
$K$ by $\beta\beta^\iota=1$, and has both torsion
frames in $V$. For the order-two generator this last
assertion needs a check: its normalized conjugator
$f_0$ is divisible by two in $V$, and
$\theta(f_0)=-f_0$ in additive notation. Taking the
harmonic primitive $g=f_0/2$ gives
$f_0=g-\theta(g)$ and an $s$-frame in $V$.
For the order-three generator, its cycle relation
is the vanishing of the order-three norm; averaging
over this odd-order group gives its frame in $V$.

Choose any lifts of these two frames to $E_2$.
Conjugating $s$ and $r$ by them defines a map from
the actual source $C_2\amalg C_3$. As in B.27,
it factors through $E_2$ and acts on its central
kernel by the natural Hopf action of $\alpha_0$,
which is identity. It is therefore an automorphism
$\widetilde\alpha_2$. The two chosen frame lifts
become central in $E_{0,2}$, so this automorphism
induces identity there.

Its square is over identity on $\overline\Gamma$
and on $K_2$, hence its difference from identity
is a continuous homomorphism
$\overline\Gamma\to K_2$. Since its image is also
trivial in $E_{0,2}$, it takes values in $L$.
But $\overline\Gamma^{\rm ab}=C_6$ and $L$ is
torsion-free, so this homomorphism is zero.
Thus $\widetilde\alpha_2^2=1$.
The possible cusp discrepancy has the form
$\widetilde\alpha_2(\tau)=\tau a$ with $a\in L$.
Squaring and using identity on $K_2$ gives $2a=0$;
again (B.30.4) gives $a=0$.
This constructs the required peripheral lift to
the complete coefficient-two centralization.

At every other coefficient prime, the normal group
$V$ is acyclic, and exactly the Cartesian-square
argument of B.27 applies, with the retained frame
pairs $(b_s,1),(b_r,1)$. Assemble these maps over
$\overline\Gamma$ and apply the modular central
lifting criterion. It gives the stated full gentle
lift, with $\tau$ fixed after the same inner
normalization as in B.27. Finite products and
compactness combine it with all the already proved
$p\geq5$ signs. The torsion exclusion is B.13,
and the failed beta cocycle uses a prime-to-two
witness. Take $\eta$ of order five and evaluate at
$(X,Y,W)=(\omega,\omega,\eta)$: the two sides are
$-1$ and $1$. Here $W=-1$ would not be a valid witness,
because the idempotent ignores the pro-two character
coordinate. The executed certificate
\texttt{pbc\_cubic2\_cocycle.py} computes the two
products in $(\mathbb Z/4)[(\mathbb Z/15)^3]$. Its
idempotent has coefficients $2$ when $i+j=0$ modulo
three and $-1$ otherwise, since $225^{-1}=1$ modulo
four. It verifies idempotence and $\beta^2=1$, and
finds a nonzero cocycle difference. Independently it
evaluates at genuine roots of orders three and five
in $(\mathbb Z/4)[t]/(t^4+t+1)$, obtaining $-1$ and
$1$. Thus this failure is an exact finite coefficient-two
obstruction.
\end{proof}

